\chapter{Clause calculi}
\label{prov:ch:claus}
\label{claus:chap}

This chapter treats the presentations whose formulas are equations and the
paired binder only. The closure binder $\Q{y}$ is not a formula constructor
in them. Its work is done by the context: a context entry is a reusable
clause $[A;\,y.\,B(y)]$ (\cref{base:sec:clauses}), and the elimination of
$\D{x}{y}B$ adds a clause where the kernel adds the formula $\Q{y}B(b,y)$.
\Cref{claus:sec:qfree} relates sequents of the native scalar core $\Nat$ to
clause sequents. \Cref{claus:sec:five} gives the five-rule calculus and its
equivalence with $\Nat$. The later sections vary one component at a time: the
arithmetic interface (\cref{claus:sec:arith}), the introduction rule
(\cref{claus:sec:intro}), the form of a context entry
(\cref{claus:sec:context}), and the number of rule families
(\cref{claus:sec:fusion}). \Cref{claus:sec:rigidity} records a separate
arithmetic observation of one of the packages, and
\cref{claus:sec:checkers} lists the checker sizes.

Every calculus of the chapter has the terms of \cref{base:sec:terms}, the
designated contradiction $0=\Suc 0$ and the ordinary reading of
\cref{base:sec:reading}. Throughout, $\equiv$ is literal identity of syntax
up to the names of bound variables. A matrix $B(x,y)$ displays the two
variables that an enclosing $\D{x}{y}$ binds, and $B(t,r)$ is the result of
substituting $t$ for $x$ and $r$ for $y$. Substitution avoids capture. A
\emph{literal equation} is a formula whose root is an equation; an equation
below a binder $D$ is not one. The numeral $\Suc 0$ is also written $1$.
The rules of $\Nat$ are cited by their names in \cref{base:sec:kernel}:
\rulename{Hyp}, the two computation equations (\rulename{I0},
\rulename{IS}), \rulename{Tr}, \rulename{QE}, \rulename{DE}, the
introduction by cases (\rulename{DI}) and the introduction of $Q$ by
progress (\rulename{QI}).

\section{From \texorpdfstring{$Q$}{Q} formulas to clauses}
\label{claus:sec:qfree}

This section draws on \pkg{raw\_clausal}, \pkg{clausal\_two\_case},
\pkg{reduced\_clause\_calculus} and \pkg{clause\_guards}.

\subsection{The \texorpdfstring{$Q$}{Q}-free form of a formula}

A native formula is \emph{$Q$-free} if it is built from equations by
$\D{x}{y}$ alone. The $Q$-free formulas are the formulas of every calculus
of this chapter.

\begin{definition}[$Q$-free form]
\label{claus:def:tau}
The translation $\tau$ from native formulas to $Q$-free formulas is
\[
 \tau(s=t)\defeq(s=t),\qquad
 \tau(\D{x}{y}B)\defeq\D{x}{y}\tau(B),\qquad
 \tau(\Q{y}B)\defeq\D{w}{y}\tau(B),
\]
where each inserted $w$ is fresh and does not occur in $B$. Terms are
unchanged. The inclusion of $Q$-free formulas into native formulas is
written $j$ and is usually left implicit. For a matrix $B(y)$ and a fresh
$w$ the \emph{closure abbreviation} is
\[
 q_yB(y)\defeq\D{w}{y}B(y),\qquad w\notin\FV(B).
\]
\end{definition}

The translation preserves free variables, commutes with substitution up to
the names of bound variables, sends equations to equations and all other
formulas to formulas with root $D$, and satisfies $\tau jF\equiv F$ for
every $Q$-free $F$. The ordinary reading of $q_yB$ is
$\exists w\,(B(0)\to\forall y\,B(y))$, which is equivalent to
$B(0)\to\forall y\,B(y)$ because the domain is inhabited by $0$. In the Lean
development $\tau$ is \Lean{unQ}.

\begin{lemma}[The round trip in $\Nat$]
\label{claus:lem:unq}
The following hold in $\Nat$, in every context.
\begin{enumerate}[label=(\roman*)]
\item $\Q{y}B\dashv\proves\D{w}{y}B$ for $w\notin\FV(B)$.
\item If $\Gamma,B\proves B'$ and $\Gamma,B'\proves B$, where the variables
$x,y$ are not free in $\Gamma$, then
$\Gamma,\D{x}{y}B\proves\D{x}{y}B'$ and $\Gamma,\Q{y}B\proves\Q{y}B'$, and
conversely.
\item $A\dashv\proves j\tau A$ for every native formula $A$.
\end{enumerate}
\end{lemma}
\begin{proof}
(i) From $\Q{y}B$: under the assumption $B(0)$ the rule \rulename{QE} gives $B(\Suc z)$,
so the introduction by cases at an arbitrary witness gives $\D{w}{y}B$.
From $\D{w}{y}B$: the rule \rulename{DE} opens $\Q{y}B(b,y)$, which is $\Q{y}B$ because
$w$ does not occur in $B$, and \rulename{Hyp} closes the branch.

(ii) Assume $\D{x}{y}B$. Apply \rulename{DE} at a fresh $b$; the branch has the
assumption $\Q{y}B(b,y)$. Introduce $\D{x}{y}B'$ by cases at the witness $b$.
Its branch assumes $B'(b,0)$. Substitution in the derivation of
$\Gamma,B'\proves B$ (\cref{base:sec:structural}) and cut give $B(b,0)$, \rulename{QE}
gives $B(b,\Suc z)$, and the substituted derivation of $\Gamma,B\proves B'$
gives $B'(b,\Suc z)$. For $Q$, pass from $\Q{y}B$ to $\D{w}{y}B$ by (i),
apply the case just proved with the unused variable $w$ in the place of
$x$, and return by (i). The converse statements are symmetric.

(iii) Induction on $A$. Equations are unchanged. For $A\equiv\D{x}{y}B$ use
(ii) and the induction hypothesis for $B$. For $A\equiv\Q{y}B$ use (ii) for
$Q$ and then (i).
\end{proof}

\subsection{Clause occurrences and their origins}
\label{claus:sec:origins}

A context of a clause calculus is a finite list of clause
\emph{occurrences} $[A;\,y.\,B(y)]$ with $y\notin\FV(A)$
(\cref{base:sec:clauses}). The clause is a context object and not a formula.
Its reading is $A\to\forall y\,B(y)$. A rule that uses a clause selects an
occurrence, so two occurrences with the same text are distinct. An ordinary
assumption $F$ is the constant clause
\[
 \operatorname{ord}(F)\defeq[\,0=0;\,v.\,F\,],\qquad v\notin\FV(F),
\]
and $\Gamma,F$ abbreviates $\Gamma,\operatorname{ord}(F)$. For a list
$\Delta$ of formulas, $\operatorname{ord}(\Delta)$ applies
$\operatorname{ord}$ to each member.

In a derivation that starts from ordinary assumptions a clause occurrence
arises in two ways only.

\begin{definition}[Origins]
\label{claus:def:origin}
An \emph{assumption occurrence} is an occurrence of the form
$\operatorname{ord}(G)$. An \emph{opened occurrence} is an occurrence of the
form $[B(b,0);\,y.\,B(b,y)]$ for a matrix $B(x,y)$ and a variable $b$. An
\emph{annotated} context is a context in which each occurrence is marked as
one of the two and has the corresponding form. The \emph{native image}
$\Gamma^{\natural}$ of an annotated context replaces each assumption
occurrence $\operatorname{ord}(G)$ by the formula $G$ and each opened
occurrence $[B(b,0);\,y.\,B(b,y)]$ by the formula $\Q{y}B(b,y)$.
\end{definition}

The mark belongs to the occurrence. A clause such as $[0=0;\,y.\,0=0]$ has
both forms, and either mark may be given to it. The rules of the calculi
below add an opened occurrence at each elimination and an assumption
occurrence at each introduction, so every context reachable from
$\operatorname{ord}(\Delta)$ has a canonical annotation.

\section{The five-rule calculus}
\label{claus:sec:five}

This section draws on \pkg{raw\_clausal}; the same lemmas are restated in
\pkg{rootfree\_clausal}, \pkg{reduced\_clause\_calculus},
\pkg{clausal\_two\_case} and \pkg{uniform\_clauses}.

\subsection{Rules}

The evaluation operation $\Phi(s,t,b)$ on raw terms is that of
\cref{comp:def:phi}. Its three rows are called the \emph{root row}
($b\equiv s$, value $t$), the \emph{zero row} ($b\equiv\It{0}{a}{u}{h}$,
$s\equiv a$, value $t$) and the \emph{successor row}
($b\equiv\It{\Suc n}{a}{u}{h}$, $s\equiv\It{n}{a}{u}{h}$, value
$h\subst{t}{u}$); they are disjoint, and the three cases of the rule
\rulename{Eval} that they give are \rulename{Root}, $\rulename{Eval}_0$ and
$\rulename{Eval}_{\Suc}$ of \cref{comp:def:e3}. The operation $\Phi_2$ is
$\Phi$ without the root row. Chapter~\ref{comp:ch} studies the rule with
contexts of native formulas; here it is combined with clause contexts.

\begin{definition}[The calculus $\mathsf{C}_3$]
\label{claus:def:c3}
Formulas are $Q$-free, contexts are lists of clause occurrences, and the
rules are the following five.
\[
 \infrule{}{\Gamma\proves s=s}\ \rulename{Refl}
 \qquad
 \infrule{\Gamma\proves s=t\quad\Gamma\proves b=r}
         {\Gamma\proves r=\Phi(s,t,b)}\ \rulename{Eval}
\]
\[
 \infrule{\Gamma\proves A}{\Gamma\proves B(t)}\ \rulename{Use}
 \qquad
 \infrule{\Gamma\proves\D{x}{y}B(x,y)\quad
          \Gamma,[B(b,0);\,y.\,B(b,y)]\proves F}
         {\Gamma\proves F}\ \rulename{Elim}
\]
\[
 \infrule{\Gamma,B(t,\eta)\proves B(t,\Suc z)}
         {\Gamma\proves\D{x}{y}B(x,y)}\ \rulename{Intro}
\]
\rulename{Eval} requires $\Phi(s,t,b)$ to be defined. \rulename{Use}
requires an occurrence $[A;\,y.\,B(y)]$ in $\Gamma$ and allows every term
$t$. In \rulename{Elim} the variable $b$ is fresh for $\Gamma$, $F$ and $B$.
In \rulename{Intro} the variable $z$ is fresh for $\Gamma$, $t$ and $B$, the
witness $t$ is an arbitrary term, and $\eta$ is the term
$\eta_{B,t}(z)$ of the witness-selected introduction (\cref{kern:def:W}):
it is $z$ if $B$ is a literal equation and $t\equiv 0$ (\emph{progress
mode}) and $0$ otherwise (\emph{case mode}). The calculus $\mathsf{C}$ is
$\mathsf{C}_3$ with $\Phi_2$ in the place of $\Phi$.
\end{definition}

There is no hypothesis rule, no computation axiom, no transport rule, no
induction rule and no structural rule. The three rows of $\Phi$ and the
mode test of \rulename{Intro} are part of the specification. The calculus
$\mathsf{C}$ is \Lean{CPrf} in the Lean development.

\subsection{Derived rules}

The lemmas of this subsection hold in $\mathsf{C}_3$ and in $\mathsf{C}$
under arbitrary clause contexts.

\begin{lemma}[Structural operations]
\label{claus:lem:struct}
An ordinary assumption can be used: $\Gamma,F\proves F$. Weakening by clause
occurrences and injective renaming of free variables are admissible. Cut for
ordinary assumptions is admissible: if $\Gamma\proves F$ and
$\Gamma,\operatorname{ord}(F),\Gamma'\proves A$ then
$\Gamma,\Gamma'\proves A$.
\end{lemma}
\begin{proof}
\rulename{Refl} proves $0=0$, and \rulename{Use} on $[0=0;\,v.\,F]$ at any
term proves $F$. Weakening and injective renaming are inductions on the
derivation after the eigenvariables are made fresh; an injective renaming
preserves the root of a matrix and the test $t\equiv 0$, so it preserves the
mode of \rulename{Intro}, and it preserves the literal matches of $\Phi$.

For cut, induct on the derivation that has the distinguished occurrence
$\operatorname{ord}(F)$. At a \rulename{Use} that selects this occurrence,
replace the node by the given derivation of $F$, weakened to the local
context; the derivation of its premise $0=0$ is discarded. At every other
node transform the premises and reapply the rule. The eigenvariables are
made fresh first. No substitution of a term in a derivation is performed.
\end{proof}

\begin{lemma}[Equality and computation]
\label{claus:lem:eq}
\rulename{Refl} and the zero and successor rows of \rulename{Eval} derive,
over every clause context, symmetry and transitivity of equality, the
inference \rulename{Root}, the two computation equations
\[
 \It{0}{a}{u}{h}=a,\qquad
 \It{\Suc n}{a}{u}{h}=h\subst{\It{n}{a}{u}{h}}{u},
\]
congruence in every term template, and transport in a term and in a
literal equation:
\[
 \infrule{s=t}{h\subst{s}{u}=h\subst{t}{u}}
 \qquad
 \infrule{s=t\quad h\subst{s}{u}=r}{r=h\subst{t}{u}}
 \qquad
 \infrule{s=t\quad a\subst{s}{u}=b\subst{s}{u}}{a\subst{t}{u}=b\subst{t}{u}}
\]
with no restriction on $s$.
\end{lemma}
\begin{proof}
The derivations of \cref{comp:sec:evalrecover}
(\cref{comp:lem:frame,comp:lem:evalalg,comp:lem:evalcomp,comp:thm:evalcong,comp:cor:evaltr})
consist of \rulename{Refl}, $\rulename{Eval}_0$, $\rulename{Eval}_{\Suc}$
and the given premises. They contain no rule that reads the context, so
they are derivations of $\mathsf{C}$ and of $\mathsf{C}_3$ over every
clause context. Transport in an equation is the one-term transport with the
template $a$, which gives $b\subst{s}{u}=a\subst{t}{u}$, followed by the
one-term transport with the template $b$.
\end{proof}

\begin{lemma}[All introductions by cases]
\label{claus:lem:cases}
For every matrix $B$ and every witness $t$ the rule
\[
 \infrule{\Gamma,B(t,0)\proves B(t,\Suc z)}{\Gamma\proves\D{x}{y}B(x,y)}
 \qquad(z\text{ fresh})
\]
is admissible.
\end{lemma}
\begin{proof}
\rulename{Intro} has this form unless $B$ is a literal equation and
$t\equiv 0$. That case is repaired as in the proof of
\cref{kern:thm:witness}, at the second name $c\equiv\It{0}{0}{u}{0}$ of
zero (\cref{kern:lem:alias}): in the context $\Gamma,B(c,0)$ transport along
$c=0$ (\cref{claus:lem:eq}) gives $B(0,0)$, the cut of
\cref{claus:lem:struct} with the given derivation gives $B(0,\Suc z)$,
transport along $0=c$ gives $B(c,\Suc z)$, and \rulename{Intro} at the
witness $c$ is in case mode. The cut is the cut for a constant clause; only
transport in an equation is used.
\end{proof}

\begin{lemma}[Transport in every formula]
\label{claus:lem:ftransport}
For every $Q$-free formula $F$ with a distinguished variable $v$, from
$\Gamma\proves s=t$ and $\Gamma\proves F\subst{s}{v}$ one derives
$\Gamma\proves F\subst{t}{v}$.
\end{lemma}
\begin{proof}
Induction on $F$. Equations are \cref{claus:lem:eq}. Let
$F\equiv\D{x}{y}B(v,x,y)$. Apply \rulename{Elim} to $F\subst{s}{v}$ at a
fresh $b$, which adds $[B(s,b,0);\,y.\,B(s,b,y)]$. Introduce
$F\subst{t}{v}$ by \cref{claus:lem:cases} at the witness $b$. Its branch
assumes $B(t,b,0)$. The induction hypothesis along $t=s$ gives $B(s,b,0)$,
\rulename{Use} at $\Suc z$ gives $B(s,b,\Suc z)$, and the induction
hypothesis along $s=t$ gives $B(t,b,\Suc z)$. Each use of the induction
hypothesis concerns the matrix $B$.
\end{proof}

\begin{lemma}[The three closure interfaces]
\label{claus:lem:closure}
With $q_yB\equiv\D{w}{y}B$, $w$ fresh, the following are derivable.
\begin{enumerate}[label=(\roman*)]
\item From $\Gamma,B(0)\proves B(\Suc z)$, $z$ fresh, derive
$\Gamma\proves q_yB$.
\item From $\Gamma\proves q_yB$ and $\Gamma\proves B(0)$ derive
$\Gamma\proves B(t)$ for every term $t$.
\item If $\Gamma$ contains an occurrence $[B(0);\,y.\,B(y)]$ then
$\Gamma\proves q_yB$.
\item For a literal equation $E(y)$: from $\Gamma,E(z)\proves E(\Suc z)$,
$z$ fresh, derive $\Gamma\proves q_yE$.
\end{enumerate}
\end{lemma}
\begin{proof}
(i) is \cref{claus:lem:cases} with the unused variable $w$ in the place of
$x$. For (ii) apply \rulename{Elim} to $q_yB$; the fresh witness does not
occur in $B$, so the added clause is $[B(0);\,y.\,B(y)]$, and \rulename{Use}
with the given derivation of $B(0)$ proves $B(t)$. For (iii), under the
assumption $B(0)$ use the retained clause at $\Suc z$ and apply (i). For
(iv) apply \rulename{Intro} to the matrix $E(y)$, in which $w$ does not
occur, at the witness $0$; the matrix is a literal equation and the witness
is literally $0$, so the rule is in progress mode.
\end{proof}

\begin{corollary}[Equation induction and numerical cases]
\label{claus:cor:induction}
For a literal equation $E(y)$, from $\Gamma\proves E(0)$ and
$\Gamma,E(z)\proves E(\Suc z)$ one derives $\Gamma\proves E(t)$. For an
arbitrary formula $F(y)$, from $\Gamma\proves F(0)$ and
$\Gamma\proves F(\Suc z)$, $z$ fresh, one derives $\Gamma\proves F(t)$. For
a literal equation $E(x,y)$ and a term $t$, from
$\Gamma,E(t,z)\proves E(t,\Suc z)$ one derives $\Gamma\proves\D{x}{y}E$.
\end{corollary}
\begin{proof}
The first two statements are (iv) and (ii), respectively (i) and (ii), of
\cref{claus:lem:closure}. For the third, (iv) gives $q_yE(t,y)$, and
\cref{claus:lem:cases} at the witness $t$ gives $\D{x}{y}E$, since under the
assumption $E(t,0)$ interface (ii) gives $E(t,\Suc z)$.
\end{proof}

Progress is available for literal equations only. Extending the progress
mode to matrices with root $D$ would give induction for arbitrary formulas
by the same argument and would be a different theory.

\subsection{Equivalence with the native scalar core}

\begin{theorem}[The five-rule calculus and $\Nat$]
\label{claus:thm:five}
Let $\mathsf{X}$ be $\mathsf{C}_3$ or $\mathsf{C}$.
\begin{enumerate}[label=(\roman*)]
\item If $\Gamma$ is an annotated clause context and
$\Gamma\proves_{\mathsf{X}}F$, then $\Gamma^{\natural}\proves_{\Nat}F$.
\item If $\Gamma$ is a native context and $\Gamma\proves_{\Nat}A$, then
$\operatorname{ord}(\tau\Gamma)\proves_{\mathsf{X}}\tau A$.
\item For $Q$-free $\Delta$ and $F$:
$\operatorname{ord}(\Delta)\proves_{\mathsf{X}}F$ if and only if
$\Delta\proves_{\Nat}F$.
\item For native $\Gamma$ and $A$:
$\operatorname{ord}(\tau\Gamma)\proves_{\mathsf{X}}\tau A$ if and only if
$\Gamma\proves_{\Nat}A$.
\end{enumerate}
All translations are effective on finite derivations and are the identity
on terms.
\end{theorem}
\begin{proof}
(i) Induction on the derivation. $\Nat$ derives reflexivity, symmetry,
transitivity and Leibniz replacement (\cref{base:sec:structural}).
\rulename{Refl} is reflexivity. The three rows of \rulename{Eval} are
derived in $\Nat$ from \rulename{I0}, \rulename{IS} and \rulename{Tr} as in
the proof of \cref{comp:thm:eval}. A \rulename{Use} of an
assumption occurrence $\operatorname{ord}(G)$ is \rulename{Hyp}. A \rulename{Use} of
an opened occurrence $[B(b,0);\,y.\,B(b,y)]$ is \rulename{QE} applied to the
assumption $\Q{y}B(b,y)$ and the translated derivation of $B(b,0)$.
\rulename{Elim} is \rulename{DE}, which adds $\Q{y}B(b,y)$, the image of the opened
occurrence. \rulename{Intro} in case mode is the introduction by cases,
with the new occurrence marked as an assumption. In progress mode the
matrix is a literal equation $E(x,y)$ and $t\equiv 0$; the translated
premise is $\Gamma^{\natural},E(0,z)\proves E(0,\Suc z)$, the introduction
of $Q$ by progress gives $\Q{y}E(0,y)$, and the introduction by cases at
the witness $0$ gives $\D{x}{y}E$, because under the assumption $E(0,0)$
the rule \rulename{QE} gives $E(0,\Suc z')$.

(ii) Induction on the derivation in $\Nat$. \rulename{Hyp} becomes \rulename{Refl} at
$0$ followed by \rulename{Use} on the constant clause. The computation
equations and \rulename{Tr} are \cref{claus:lem:eq}.
\rulename{QE} is interface (ii) of \cref{claus:lem:closure}, since
$\tau(\Q{y}B)\equiv q_y\tau(B)$ and $\tau$ commutes with substitution. The
introduction by cases is \cref{claus:lem:cases}. The introduction of
$\Q{y}E$ by progress is interface (iv), whose conclusion is
$q_yE\equiv\tau(\Q{y}E)$. For \rulename{DE} with major premise $\D{x}{y}B$ and branch
$\Gamma,\Q{y}B(b,y)\proves A$, translate the major premise and apply
\rulename{Elim}, which adds $[\tau B(b,0);\,y.\,\tau B(b,y)]$. In the
extended context interface (iii) derives $q_y\tau B(b,y)$, which is
$\tau(\Q{y}B(b,y))$. Weaken the translated branch to the extended context
and cut this derivation in for its ordinary assumption
(\cref{claus:lem:struct}). \rulename{Elim} then concludes $\tau A$.

(iii) From left to right use (i) with the canonical annotation, for which
$\operatorname{ord}(\Delta)^{\natural}=\Delta$. From right to left use (ii);
$\tau$ is the identity on $Q$-free formulas.

(iv) From right to left is (ii). From left to right, (i) gives
$\tau\Gamma\proves_{\Nat}\tau A$, and \cref{claus:lem:unq}(iii) with cut in
$\Nat$ replaces $\tau\Gamma$ by $\Gamma$ and $\tau A$ by $A$.
\end{proof}

\begin{corollary}[Both formula composites]
\label{claus:cor:composites}
The composite $\tau j$ is the identity on $Q$-free formulas, and
$A\dashv\proves_{\Nat}j\tau A$ for every native formula $A$. Both term
translations are identities. In particular
$\proves_{\mathsf{C}}0=\Suc 0$ if and only if $\proves_{\Nat}0=\Suc 0$.
\end{corollary}

\begin{remark}
\label{claus:rem:fivepkgs}
In the Lean development (iv) is \Lean{cprf\_iff\_nprf} and (iii) is
\Lean{cprf\_iff\_nprf\_of\_qfree}, for the calculus $\mathsf{C}$
(\Lean{NativeClausal}). There the clause calculus is defined over all
native formulas, and \Lean{CPrf.qfree} states that it derives only
$Q$-free formulas from $Q$-free clauses. \pkg{raw\_clausal} states the
theorem for $\mathsf{C}_3$ against the six-rule calculus with
\rulename{Refl} and \rulename{Eval} of \pkg{computational\_equality};
\pkg{clausal\_two\_case} and \pkg{uniform\_clauses} state it for
$\mathsf{C}$ against the same calculus with two evaluation rows, and
\pkg{clause\_guards} against the six-rule calculus with the diagonal
introduction. The proof above compares with $\Nat$ directly; the cases are
the same. The Lean proof of completeness translates the six-rule calculus
rule by rule.
\end{remark}

\begin{proposition}[Substitution]
\label{claus:prop:subst}
If $\Gamma\proves_{\mathsf{X}}F$ and $\sigma$ is a substitution of terms for
free variables, then $\Gamma\sigma\proves_{\mathsf{X}}F\sigma$, for
$\mathsf{X}$ either of $\mathsf{C}_3$, $\mathsf{C}$ and every clause context
$\Gamma$. The rule at a node of the derivation is not always preserved.
\end{proposition}
\begin{proof}
Induction on the derivation, with eigenvariables made fresh for $\sigma$.
\rulename{Refl} is immediate, and
$\Phi(s\sigma,t\sigma,b\sigma)\equiv\Phi(s,t,b)\sigma$ by the same row
(\cref{comp:prop:phisubst}). \rulename{Use} and \rulename{Elim} commute
with $\sigma$. An
\rulename{Intro} in progress mode keeps its equation root and its literal
witness $0$. An \rulename{Intro} in case mode whose matrix has root $D$, or
whose witness $t$ has $t\sigma\not\equiv 0$, remains in case mode. In the
remaining case the matrix is a literal equation, $t\not\equiv 0$ and
$t\sigma\equiv 0$; the substituted premise is the premise of an
introduction by cases at the witness $0$, and \cref{claus:lem:cases}
applies.
\end{proof}

A substitution can thus change the mode of an introduction, and the
derivation after substitution may contain the nodes of
\cref{claus:lem:cases} in the place of one \rulename{Intro}.
\Cref{claus:thm:ksubst} gives a calculus in which the tree is preserved.

\begin{remark}[Clauses that are not of the two origins]
\label{claus:rem:external}
\Cref{claus:thm:five}(i) concerns annotated contexts. An arbitrary clause
$[A;\,y.\,B]$ supplied from outside is represented in $\Nat$ by the derived
implication and universal quantifier of the native logical interface
(\cref{base:thm:bridge}), as the formula that reads $A\to\forall y\,B$, and
\rulename{Use} becomes the elimination of the implication followed by the
elimination of the universal quantifier. The equivalences (iii) and (iv)
do not need such clauses. They cover every sequent with ordinary
assumptions and not only closed theorems.
\end{remark}

\section{Variants of the arithmetic interface}
\label{claus:sec:arith}

This section draws on \pkg{rootfree\_clausal},
\pkg{reduced\_clause\_calculus}, \pkg{clausal\_two\_case},
\pkg{clausal\_single\_evaluation}, \pkg{raw\_clausal},
\pkg{uniform\_clauses} and \pkg{clause\_contexts}; each subsection names
its sources. The variants are: an evaluation rule with two rows, an
evaluation rule with one scheme on zero-normal terms, and reflexivity
stated as a clause, on raw and on zero-normal terms.

\subsection{Two evaluation cases}
\label{claus:sec:twocase}

This subsection draws on \pkg{rootfree\_clausal},
\pkg{reduced\_clause\_calculus} and \pkg{clausal\_two\_case}, which present
the same calculus $\mathsf{C}$. Its two evaluation rows are the rules
$\rulename{Eval}_0$ and $\rulename{Eval}_{\Suc}$ of \cref{comp:def:e3}. The
root row is derived by the identity frame of \cref{comp:lem:frame}: with
$K(s)\equiv\It{0}{s}{u}{0}$, the derivation
\[
 \mathrm{RC}(p,q)\defeq
 \mathrm{Ev}_0\bigl(p,\ \mathrm{Ev}_0(q,\ \mathrm{R}_{K(s)})\bigr)
\]
of \cref{comp:thm:root} turns $p:s=t$ and $q:s=r$ into $r=t$ by two
instances of $\rulename{Eval}_0$ and one \rulename{Refl} leaf. It contains
no rule that reads the context, so it is a derivation of $\mathsf{C}$ over
every clause context. The body of $K(s)$ is never evaluated; any body
serves.

\begin{theorem}[Deletion of the root row]
\label{claus:thm:twocase}
For every clause context $\Gamma$ and every formula $F$,
\[
 \Gamma\proves_{\mathsf{C}_3}F\quad\Longleftrightarrow\quad
 \Gamma\proves_{\mathsf{C}}F .
\]
The translation from $\mathsf{C}_3$ to $\mathsf{C}$ leaves the end sequent
unchanged and turns a derivation with $N$ nodes, $k$ of them uses of the
root row, into a derivation with exactly $N+2k\le 3N$ nodes. The converse
translation is the inclusion.
\end{theorem}
\begin{proof}
Every rule of $\mathsf{C}$ is a rule of $\mathsf{C}_3$. Conversely replace
each use of the root row by $\mathrm{RC}$ and copy every other node, as in
\cref{comp:thm:root}; the count is the count of that theorem. What the
clause setting adds is that the replacement is made under arbitrary
retained clauses, also those opened inside \rulename{Elim} and
\rulename{Intro}: no variable is bound or freed, and the matrices and
witnesses of \rulename{Intro} are unchanged, so its modes and all
freshness conditions are unchanged.
\end{proof}

The count is of inference nodes and bounds no serialized size
(\cref{comp:thm:root}).

\begin{proposition}[Three new nodes are necessary]
\label{claus:prop:threenodes}
Let $x,y,z$ be distinct variables, and consider the fragment with the rules
\rulename{Refl}, $\rulename{Eval}_0$ and $\rulename{Eval}_{\Suc}$ and with
two given leaves $p:x=y$ and $q:x=z$. Every derivation of $z=y$ in this
fragment contains at least two \rulename{Eval} nodes and one
\rulename{Refl} leaf. $\mathrm{RC}(p,q)$ attains the bound.
\end{proposition}
\begin{proof}
The conclusion $z=y$ is not $p$, not $q$ and not reflexive, so the last
inference is an \rulename{Eval}. In both rows the second premise has the
form $b=r$ with $b$ an iterator, and the conclusion has $r$ as its left
side. The second premise of the last inference is therefore $b=z$ with $b$
an iterator. It is not $p$ or $q$, whose left side is the variable $x$, and
it is not a \rulename{Refl} leaf, since an iterator is not the variable
$z$. So it ends in another \rulename{Eval}. Follow second premises upwards
from there. Each has an iterator as its left side, so the path does not end
at $p$ or $q$; it ends at a \rulename{Refl} leaf.
\end{proof}

The bound concerns this fragment and this cost convention. It does not
bound derivations that use the logical rules or previously derived lemmas,
and it is not a bound on the size of a checker.

\begin{proposition}[Neither remaining row is redundant]
\label{claus:prop:rows}
In $\mathsf{C}$ without $\rulename{Eval}_0$ the equation
$\It{0}{0}{u}{0}=0$ is not derivable. In $\mathsf{C}$ without
$\rulename{Eval}_{\Suc}$ the equation
$\It{\Suc 0}{0}{u}{\Suc u}=\Suc 0$ is not derivable.
\end{proposition}
\begin{proof}
The two interpretations of \cref{comp:prop:evalsharp} apply, with a clause
read as $A\to\forall y\,B(y)$: under any compositional interpretation of
the iterator in the standard natural numbers \rulename{Refl},
\rulename{Use}, \rulename{Elim} and both modes of \rulename{Intro} are
sound, each interpretation makes the remaining row sound, and the displayed
equations are false, the second because its left side has the value $0$.
\end{proof}

\begin{remark}
\label{claus:rem:twopkgs}
The three packages state \cref{claus:thm:twocase} with the
frame $\It{0}{s}{u}{0}$; \pkg{uniform\_clauses}, \pkg{clause\_guards},
\pkg{one\_body} and \pkg{prefix\_clauses} use $\It{0}{s}{u}{u}$. The frame
argument is that of \pkg{equality\_from\_iteration}
(\cref{comp:lem:frame}), applied inside the
clause calculus. \Cref{claus:prop:threenodes} is in
\pkg{reduced\_clause\_calculus} and, in outline, in
\pkg{clausal\_two\_case}; \cref{claus:prop:rows} is in
\pkg{clausal\_two\_case}. \pkg{rootfree\_clausal} also gives the direct
proof of \cref{claus:prop:subst}.
\end{remark}

From here on $\Phi_2$ has two rows, and the lemmas of \cref{claus:sec:five}
are used for $\mathsf{C}$.

\subsection{Zero-normal terms and one evaluation scheme}
\label{claus:sec:zero}

This subsection draws on \pkg{clausal\_single\_evaluation}, on the
companion calculi of \pkg{reduced\_clause\_calculus} and
\pkg{rootfree\_clausal}, and on \pkg{clause\_contexts}. The zero computation
$\It{0}{a}{u}{h}=a$ is moved from the rules into the syntax.

Zero-normal terms, the smart constructor $\overline{I}$ and the
normalization $\nf_0$ are those of \cref{comp:def:nfzero}, and normal
substitution is the normalizing substitution
$t\langle s/x\rangle\defeq\nf_0(t\subst{s}{x})$ of \cref{comp:def:nusubst}.
The map $\nf_0$ acts on formulas, clauses and contexts through their terms,
and $B\langle t,r\rangle$ is the corresponding instance of a matrix. Two
facts are used below. Normalization is coherent with substitution,
\[
 \nf_0(t\sigma)\equiv\nf_0\bigl(\nf_0(t)\,(\nf_0\circ\sigma)\bigr)
\]
(\cref{comp:lem:nfzerosubst}), so normal substitutions compose
(\cref{comp:lem:laws}). And $\Nat$ derives $t=\nf_0(t)$ for every raw term
$t$ (\cref{comp:lem:conversion}), so that $\Gamma,F\proves_{\Nat}\nf_0(F)$
and $\Gamma,\nf_0(F)\proves_{\Nat}F$ for every formula $F$ and context
$\Gamma$ (\cref{comp:lem:nusound,comp:prop:nfzeroadm}).


\begin{definition}[The calculus $\mathsf{C}^{0}$]
\label{claus:def:c0}
Terms, formulas and clauses are zero-normal, contexts and the notion of
origin are as before, and every substitution performed by a rule is normal
substitution. The rules are
\[
 \infrule{}{\Gamma\proves s=s}\ \rulename{Refl}
 \qquad
 \infrule{\Gamma\proves A}{\Gamma\proves B\langle t/y\rangle}\ \rulename{Use}
\]
\[
 \infrule{\Gamma\proves\overline{I}(n,a;\,u.\,h)=t\quad
          \Gamma\proves\It{\Suc n}{a}{u}{h}=r}
         {\Gamma\proves r=h\langle t/u\rangle}\ \rulename{Eval}
\]
\[
 \infrule{\Gamma\proves\D{x}{y}B\quad
          \Gamma,[B\langle b,0\rangle;\,y.\,B\langle b,y\rangle]\proves F}
         {\Gamma\proves F}\ \rulename{Elim}
\]
\[
 \infrule{\Gamma,B\langle t,\eta\rangle\proves B\langle t,\Suc z\rangle}
         {\Gamma\proves\D{x}{y}B}\ \rulename{Intro}
\]
with the side conditions of \cref{claus:def:c3}; in particular $\eta$ is
$z$ if $B$ is a literal equation and $t\equiv 0$, and $0$ otherwise.
\end{definition}

There is one evaluation scheme, the rule \rulename{Step} of
\cref{comp:def:step}; it has no root row and no zero row. The calculus is
$\mathsf{C5}$ of \pkg{clausal\_single\_evaluation}, it is \Lean{ZCPrf} in
the Lean development, and it coincides with the calculus
$\mathsf{C}_{\mathrm{Step}}$ that \cref{comp:def:cstep} takes from
\pkg{uniform\_step}; \cref{comp:sec:clausal} derives the native rules in it
for that package and proves statement (iv) of \cref{claus:thm:zero} below
through $\Teq$.

\begin{lemma}[Derived rules of $\mathsf{C}^{0}$]
\label{claus:lem:c0}
The following hold in $\mathsf{C}^{0}$ under arbitrary clause contexts.
\begin{enumerate}[label=(\roman*)]
\item The structural operations of \cref{claus:lem:struct}.
\item From $p:s=t$ and $q:s=r$ one derives $r=t$ by two \rulename{Eval}
nodes and one \rulename{Refl} leaf. Hence symmetry and transitivity.
\item The successor computation
$\It{\Suc n}{a}{u}{h}=h\langle\overline{I}(n,a;\,u.\,h)/u\rangle$.
\item From $p:s=t$ one derives $h\langle s/u\rangle=h\langle t/u\rangle$ by
two \rulename{Eval} nodes and two \rulename{Refl} leaves, with $p$ used
once.
\item From $p:s=t$ and $q:h\langle s/u\rangle=r$ one derives
$r=h\langle t/u\rangle$ by four \rulename{Eval} nodes and three
\rulename{Refl} leaves, with $p$ and $q$ used once each. Two such steps
give transport in an equation.
\item The introduction by cases of \cref{claus:lem:cases} for every matrix
and every witness.
\item Transport in every formula (\cref{claus:lem:ftransport}), the closure
interfaces (\cref{claus:lem:closure}) and \cref{claus:cor:induction}.
\end{enumerate}
\end{lemma}
\begin{proof}
(i) The proofs of \cref{claus:lem:struct} apply; injective renaming
preserves normal terms.

(ii) to (v) The rule \rulename{Eval} of $\mathsf{C}^{0}$ is the rule
\rulename{Step} of \cref{comp:def:step}, and the derivations of
\cref{comp:lem:stepderived} consist of \rulename{Refl}, \rulename{Step}
and the given premises, so they are derivations of $\mathsf{C}^{0}$ over
every clause context: (ii) is the pivot with the frame
$\It{\Suc 0}{s}{u}{u}$, with symmetry and transitivity; (iii) is the
successor computation; (iv) is the congruence; (v) is the transport.
Replacing one node of a one-term transport rule by the derivation of (v)
adds six nodes and duplicates neither premise.

(vi) Only a literal equation at the witness $0$ needs an argument. Let
$Z\equiv\It{\Suc 0}{0}{u}{0}$. It is zero-normal and is not literally $0$
(\cref{comp:lem:alias}), and \rulename{Eval} on the reflexivities of $0$
and of $Z$ proves $Z=0$. The proof of \cref{claus:lem:cases} applies with
$Z$ in the place of $c$. The raw term $\It{0}{0}{u}{0}$ cannot be used: it
normalizes to $0$.

(vii) The proofs of \cref{claus:lem:ftransport,claus:lem:closure} use only
the items above, \rulename{Use}, \rulename{Elim} and \rulename{Intro}.
\end{proof}


\begin{theorem}[The zero-normal calculus and $\Nat$]
\label{claus:thm:zero}
\begin{enumerate}[label=(\roman*)]
\item For every raw clause context $\Gamma$ and raw formula $F$: if
$\Gamma\proves_{\mathsf{C}}F$ then
$\nf_0(\Gamma)\proves_{\mathsf{C}^{0}}\nf_0(F)$.
\item If $\Gamma$ is an annotated zero-normal context and
$\Gamma\proves_{\mathsf{C}^{0}}F$, then $\Gamma^{\natural}\proves_{\Nat}F$.
Here an opened occurrence is
$[B\langle b,0\rangle;\,y.\,B\langle b,y\rangle]$ and its native image is
$\Q{y}B(b,y)$.
\item For raw $Q$-free $\Delta$ and $F$:
$\operatorname{ord}(\nf_0\Delta)\proves_{\mathsf{C}^{0}}\nf_0(F)$ if and only
if $\Delta\proves_{\Nat}F$. In particular, for zero-normal $\Delta$ and
$F$, $\operatorname{ord}(\Delta)\proves_{\mathsf{C}^{0}}F$ if and only if
$\Delta\proves_{\Nat}F$.
\item For native $\Gamma$ and $A$:
$\operatorname{ord}(\nf_0\tau\Gamma)\proves_{\mathsf{C}^{0}}\nf_0(\tau A)$ if
and only if $\Gamma\proves_{\Nat}A$.
\end{enumerate}
\end{theorem}
\begin{proof}
(i) Induction on the derivation in $\mathsf{C}$, normalizing every sequent.
\rulename{Refl} is \rulename{Refl}. For $\rulename{Eval}_0$ the normalized
premises are $\nf_0(a)=\nf_0(t)$ and $\nf_0(a)=\nf_0(r)$, since
$\nf_0(\It{0}{a}{u}{h})\equiv\nf_0(a)$, and \cref{claus:lem:c0}(ii) gives
$\nf_0(r)=\nf_0(t)$. For $\rulename{Eval}_{\Suc}$ the normalized premises are
\[
 \overline{I}(\nf_0 n,\nf_0 a;\,u.\,\nf_0 h)=\nf_0(t),\qquad
 \It{\Suc\nf_0 n}{\nf_0 a}{u}{\nf_0 h}=\nf_0(r),
\]
because the counter $\Suc\nf_0(n)$ is not $0$. \rulename{Eval} of
$\mathsf{C}^{0}$ gives $\nf_0(r)=\nf_0(h)\langle\nf_0(t)/u\rangle$, which is
$\nf_0(r)=\nf_0(h\subst{t}{u})$ by \cref{comp:lem:nfzerosubst}. A \rulename{Use} of
$[A;\,y.\,B]$ at $t$ becomes a \rulename{Use} of $[\nf_0 A;\,y.\,\nf_0 B]$ at
$\nf_0(t)$, with conclusion
$\nf_0(B)\langle\nf_0(t)/y\rangle\equiv\nf_0(B\subst{t}{y})$. For \rulename{Elim}
the normal form of the added clause $[B(b,0);\,y.\,B(b,y)]$ is
$[\nf_0(B)\langle b,0\rangle;\,y.\,\nf_0(B)\langle b,y\rangle]$, the clause that
\rulename{Elim} of $\mathsf{C}^{0}$ adds for $\D{x}{y}\nf_0(B)$. For
\rulename{Intro} with matrix $B$ and witness $t$ the normalized premise is
$\nf_0(\Gamma),\nf_0(B)\langle\nf_0 t,\eta\rangle\proves
\nf_0(B)\langle\nf_0 t,\Suc z\rangle$. The matrix $\nf_0(B)$ is a literal
equation exactly when $B$ is one. If the mode of \rulename{Intro} for
$(\nf_0 B,\nf_0 t)$ is the mode for $(B,t)$, apply \rulename{Intro} of
$\mathsf{C}^{0}$. The modes differ only when $B$ is a literal equation,
$t\not\equiv 0$ and $\nf_0(t)\equiv 0$. Then the source is in case mode, the
normalized premise is the premise of an introduction by cases at the
witness $0$, and \cref{claus:lem:c0}(vi) applies.

(ii) Induction on the derivation, as in \cref{claus:thm:five}(i).
\rulename{Refl} is reflexivity. For \rulename{Eval}, $\Nat$ derives
$\It{n}{a}{u}{h}=\overline{I}(n,a;\,u.\,h)$ by the zero computation or by
reflexivity, hence $\It{n}{a}{u}{h}=t$ from the first premise; the
successor computation, Leibniz replacement and \cref{comp:lem:nusound} give
$\It{\Suc n}{a}{u}{h}=h\subst{t}{u}=\nf_0(h\subst{t}{u})$, and the second
premise gives $r=h\langle t/u\rangle$. A \rulename{Use} of an assumption
occurrence is \rulename{Hyp}. For a \rulename{Use} of an opened occurrence the
translated premise is $B\langle b,0\rangle$; \cref{comp:lem:nusound} gives
$B(b,0)$, \rulename{QE} on $\Q{y}B(b,y)$ gives $B(b,t)$, and \cref{comp:lem:nusound}
gives $B\langle b,t\rangle$. \rulename{Elim} is \rulename{DE}. For \rulename{Intro}
the translated premise is
$\Gamma^{\natural},B\langle t,\eta\rangle\proves B\langle t,\Suc z\rangle$;
\cref{comp:lem:nusound} and cut in $\Nat$ give
$\Gamma^{\natural},B(t,\eta)\proves B(t,\Suc z)$, and the two modes are
treated as in \cref{claus:thm:five}(i).

(iii) If $\Delta\proves_{\Nat}F$ then
$\operatorname{ord}(\Delta)\proves_{\mathsf{C}}F$ by
\cref{claus:thm:five}(iii), and (i) gives
$\operatorname{ord}(\nf_0\Delta)\proves_{\mathsf{C}^{0}}\nf_0(F)$, because
$\nf_0(\operatorname{ord}(G))\equiv\operatorname{ord}(\nf_0 G)$. Conversely
(ii) gives $\nf_0(\Delta)\proves_{\Nat}\nf_0(F)$, and \cref{comp:lem:nusound}
with cut gives $\Delta\proves_{\Nat}F$.

(iv) Combine (iii) for $\tau\Gamma$ and $\tau A$ with
\cref{claus:lem:unq}(iii).
\end{proof}

\begin{corollary}[Substitution]
\label{claus:cor:zerosubst}
Normal substitution of terms for free variables is admissible in
$\mathsf{C}^{0}$ for sequents with ordinary assumptions.
\end{corollary}
\begin{proof}
By \cref{claus:thm:zero}(iii), substitution in $\Nat$
(\cref{base:sec:structural}) and
$\nf_0(F\sigma)\equiv\nf_0(F)\langle\nf_0\circ\sigma\rangle$ from
\cref{comp:lem:nfzerosubst}. A witness may become the literal $0$ under the
substitution, so the rule at a node is not always preserved.
\end{proof}

\begin{remark}
\label{claus:rem:zeropkgs}
In the Lean development (iv) is \Lean{zcprf\_iff\_nprf}, (iii) is
\Lean{zcprf\_iff\_nprf\_of\_qfree} and the translation of (i) is
\Lean{CPrf.toZC} (\Lean{NativeZeroClausal}). \pkg{clausal\_single\_evaluation}
proves (iii) for zero-normal sequents by comparison with the kernel
$\Nat^{0}$ on zero-normal terms (\cref{comp:thm:zero}); the cases are those
of \cref{claus:thm:five} with
\cref{claus:lem:c0} in the place of the lemmas of \cref{claus:sec:five}.
The proof above normalizes derivations of $\mathsf{C}$, which is the route
of the companion of \pkg{reduced\_clause\_calculus} and of the Lean proof.
The change of mode in (i) and its repair at the witness
$\It{\Suc 0}{0}{u}{0}$ are stated in \pkg{reduced\_clause\_calculus} and
\pkg{clausal\_single\_evaluation}. The count in
\cref{claus:lem:c0}(iv) and (v) is from \pkg{clausal\_single\_evaluation};
the congruence of \cref{claus:lem:eq} has four evaluations because it
joins two separately evaluated endpoints. \pkg{rootfree\_clausal} states
\cref{claus:lem:c0}(ii) to (iv) for its companion, whose introduction is
that of \cref{claus:sec:fixedzero}.
\end{remark}

\subsection{Reflexivity as one fixed clause}
\label{claus:sec:uniform}

This subsection draws on \pkg{raw\_clausal} for the obstruction and on
\pkg{uniform\_clauses} for the calculus.

\begin{theorem}[Closed seeds do not replace reflexivity]
\label{claus:thm:reflobs}
Let $S$ be a finite set of closed equations that are true in the standard
natural numbers, and let $M\subseteq\mathbb{N}$ be the finite set of values
of their sides. Let $\mathsf{C}_3^{S}$ be $\mathsf{C}_3$ with
\rulename{Refl} deleted and the members of $S$ added as axioms in every
context. If $k\notin M$, then $\mathsf{C}_3^{S}$ does not derive
$\Suc^k0=\Suc^k0$ from the empty context. In particular the calculus with
the single axiom $0=0$ in the place of \rulename{Refl} does not derive
$\Suc 0=\Suc 0$, and the calculus with \rulename{Refl} deleted derives no
equation from the empty context.
\end{theorem}
\begin{proof}
Evaluate terms in the standard natural numbers, the iterator by ordinary
iteration. Give the equation atom the truth condition
\[
 [\![s=t]\!]_\nu\quad\iffdef\quad
 [\![s]\!]_\nu=[\![t]\!]_\nu\in M .
\]
Interpret $\D{x}{y}B$ as
$\exists x\,([\![B(x,0)]\!]\to\forall y\,[\![B(x,y)]\!])$ and a clause
$[A;\,y.\,B]$ as $[\![A]\!]\to\forall y\,[\![B(y)]\!]$, with two-valued
truth. A sequent is valid if its formula is true under every assignment
that makes its clauses true. Every rule preserves validity.

Each seed is true. For \rulename{Eval} let the premises $s=t$ and $b=r$ be
true, so that $[\![s]\!]=[\![t]\!]\in M$ and $[\![b]\!]=[\![r]\!]\in M$. In
each row $\Phi(s,t,b)$ has the value of $b$: in the root row
$[\![t]\!]=[\![s]\!]=[\![b]\!]$; in the zero row
$[\![t]\!]=[\![a]\!]=[\![\It{0}{a}{u}{h}]\!]$; in the successor row
$[\![h\subst{t}{u}]\!]=[\![h\subst{s}{u}]\!]=[\![\It{\Suc n}{a}{u}{h}]\!]$,
because $s\equiv\It{n}{a}{u}{h}$ and $t$ have the same value. Hence the
conclusion $r=\Phi(s,t,b)$ is true. \rulename{Use} is the elimination of
an implication and of a universal quantifier. \rulename{Elim} is
existential elimination; the fresh variable takes the value of the
witness, and neither the context nor the conclusion depends on it. For
\rulename{Intro} fix an assignment that satisfies $\Gamma$. The clause
$\operatorname{ord}(F)$ is true whenever $F$ is.
% ed.: raw_clausal treats the local assumption ord(F) as F; when 0 is not
% in M the antecedent 0=0 is false and ord(F) is vacuously true, which the
% preceding sentence covers.
In case mode, if
$B(t,0)$ is true then the premise makes $B(t,\Suc n)$ true for every
standard $n$; every number is $0$ or a successor, so every instance is
true and the conditional closure holds at the value of $t$. If $B(t,0)$ is
false the closure holds trivially. In progress mode the premise states
that $B(0,n)$ implies $B(0,\Suc n)$ for every standard $n$; induction on
the standard numbers gives every instance from $B(0,0)$. The argument does
not use that the atoms mean identity.

So every theorem of $\mathsf{C}_3^{S}$ is true in the interpretation, and
$\Suc^k0=\Suc^k0$ is false for $k\notin M$. For the two special cases take
$S=\{0=0\}$, $M=\{0\}$, and $S=\emptyset$.
\end{proof}

The interpretation is a counter-interpretation of an inference system. It
is not a structure in which the equality symbol is identity, and it is not
the arithmetic interpretation of the calculus. For contexts of native
formulas the case $S=\emptyset$ is proved syntactically in
\cref{comp:prop:reflsharp}; the theorem adds the closed seeds. The clause
\[
 [\,0=0;\,y.\,y=y\,]
\]
is false in it. Retained in every context, together with a derivation of
$0=0$, it restores \rulename{Refl}: \rulename{Use} at $t$ gives $t=t$. This
relocates universal reflexivity into the initial context; the axiom and
the cost of validating it remain part of the specification. The calculus
that follows makes this choice and also removes the neutral antecedent.

\begin{definition}[The calculus $\mathsf{V}$]
\label{claus:def:v}
A \emph{uniform clause} is $[\Pi;\,y.\,B(y)]$ where $\Pi$ is empty or one
formula and $y$ does not occur in $\Pi$; it reads $\forall y\,B(y)$,
respectively $A\to\forall y\,B(y)$. The \emph{reflexivity clause} and the
clause of an assumption are
\[
 \mathcal{R}\defeq[\,;\,y.\,y=y\,],\qquad
 \langle A\rangle\defeq[\,;\,v.\,A\,]\quad(v\notin\FV(A)).
\]
The rules are
\[
 \infrule{\Gamma\proves A_1\ \cdots\ \Gamma\proves A_k}
         {\Gamma\proves B(t)}\ \rulename{Use}
 \qquad
 \infrule{\Gamma\proves s=t\quad\Gamma\proves b=r}
         {\Gamma\proves r=\Phi_2(s,t,b)}\ \rulename{Eval}
\]
\[
 \infrule{\Gamma\proves\D{x}{y}B\quad
          \Gamma,[B(b,0);\,y.\,B(b,y)]\proves F}
         {\Gamma\proves F}\ \rulename{Elim}
\]
\[
 \infrule{\Gamma,\langle B(t,\eta)\rangle\proves B(t,\Suc z)}
         {\Gamma\proves\D{x}{y}B}\ \rulename{Intro}
\]
where \rulename{Use} selects an occurrence $[A_1,\dots,A_k;\,y.\,B]$ of
$\Gamma$ with $k\in\{0,1\}$, and the side conditions and $\eta$ are those
of \cref{claus:def:c3}. For a context $\Gamma$ of uniform clauses,
$\Gamma^{\circ}$ replaces each occurrence $[\,;\,y.\,B]$ by
$[0=0;\,y.\,B]$.
\end{definition}

There are four rule families and one fixed axiom clause. \rulename{Use} of
$\mathcal{R}$ at $t$ gives $t=t$, and \rulename{Use} of $\langle A\rangle$
gives $A$; both have no premise. Reflexivity is not eliminated as a
principle; it is the clause $\mathcal{R}$. The calculus is \Lean{VPrf} in
the Lean development.

\begin{theorem}[Uniform clauses]
\label{claus:thm:uniform}
For every context $\Gamma$ of uniform clauses and every formula $A$,
\[
 \mathcal{R},\Gamma\proves_{\mathsf{V}}A\quad\Longleftrightarrow\quad
 \Gamma^{\circ}\proves_{\mathsf{C}}A .
\]
Hence, for $Q$-free $\Delta$ and $F$,
$\mathcal{R},\langle\Delta\rangle\proves_{\mathsf{V}}F$ if and only if
$\Delta\proves_{\Nat}F$. The translations leave terms and formulas
unchanged. The translation from $\mathsf{C}$ to $\mathsf{V}$ does not
increase the number of nodes; the converse adds at most one
\rulename{Refl} node for each use of a clause without antecedent.
\end{theorem}
\begin{proof}
The occurrences of $\Gamma$ and of $\Gamma^{\circ}$ correspond by
position, and $\mathcal{R}$ is the first occurrence on the left. Both
directions are inductions on derivations, over all $\Gamma$.

From right to left: \rulename{Refl} at $t$ becomes \rulename{Use} of
$\mathcal{R}$ at $t$. A \rulename{Use} of an occurrence $[0=0;\,y.\,B]$ that
corresponds to $[\,;\,y.\,B]$ becomes the \rulename{Use} without premise,
and the derivation of $0=0$ is discarded. A \rulename{Use} of an occurrence
that corresponds to a clause with an antecedent keeps its translated
premise. \rulename{Eval} is copied. \rulename{Elim} adds the same clause on
both sides, and \rulename{Intro} adds $\operatorname{ord}(B(t,\eta))$,
which is $\langle B(t,\eta)\rangle^{\circ}$. Witnesses, modes and
eigenvariables are unchanged.

From left to right: a \rulename{Use} of $\mathcal{R}$ at $t$ becomes
\rulename{Refl} at $t$. A \rulename{Use} of $[\,;\,y.\,B]$ becomes a
\rulename{Use} of $[0=0;\,y.\,B]$ with premise \rulename{Refl} at $0$. The
other rules are copied.

The second statement follows from \cref{claus:thm:five}(iii), since
$\langle\Delta\rangle^{\circ}=\operatorname{ord}(\Delta)$.
\end{proof}

\begin{remark}
\label{claus:rem:uniformpkgs}
In the Lean development the first statement is \Lean{vprf\_iff\_cprf} and
the second is \Lean{vprf\_iff\_nprf\_of\_qfree}
(\Lean{NativeUniformClauses}). In the native reading of a context of
$\mathsf{V}$ that is reachable from assumptions, $\mathcal{R}$ supplies
reflexivity, $\langle A\rangle$ is the assumption $A$ and an opened clause
is $\Q{y}B(b,y)$, so \rulename{Use} is reflexivity, \rulename{Hyp} or \rulename{QE}.
\end{remark}

\subsection{Four schemes on zero-normal terms}
\label{claus:sec:four}

This subsection draws on \pkg{clause\_contexts}.

\begin{definition}[The calculus $\mathsf{V}^{0}$]
\label{claus:def:v0}
$\mathsf{V}^{0}$ has zero-normal terms, normal substitution, uniform
clauses and the clause $\mathcal{R}$, with the four schemes
\rulename{Use} (as in $\mathsf{V}$), \rulename{Step} (the rule
\rulename{Eval} of $\mathsf{C}^{0}$), \rulename{Open} (the rule
\rulename{Elim} of $\mathsf{C}^{0}$) and \rulename{Close} (the rule
\rulename{Intro} of $\mathsf{C}^{0}$ with the assumption stored as
$\langle B\langle t,\eta\rangle\rangle$).
\end{definition}

\begin{theorem}[Clause factorization]
\label{claus:thm:four}
For every zero-normal context $\Gamma$ of uniform clauses and every
zero-normal formula $A$,
$\mathcal{R},\Gamma\proves_{\mathsf{V}^{0}}A$ if and only if
$\Gamma^{\circ}\proves_{\mathsf{C}^{0}}A$. Hence, for raw $Q$-free $\Delta$
and $F$,
$\mathcal{R},\langle\nf_0\Delta\rangle\proves_{\mathsf{V}^{0}}\nf_0(F)$ if and
only if $\Delta\proves_{\Nat}F$.
\end{theorem}
\begin{proof}
The proof of \cref{claus:thm:uniform} applies; it does not inspect the
evaluation rule or the terms. The second statement is
\cref{claus:thm:zero}(iii).
\end{proof}

In the Lean development $\mathsf{V}^{0}$ is \Lean{ZVPrf} and the two
statements are \Lean{zvprf\_iff\_zcprf} and \Lean{zvprf\_iff\_nprf}.

\begin{remark}[The atomic case without cut]
\label{claus:rem:nocut}
\pkg{clause\_contexts} recovers the introduction by cases of a literal
equation $E(x,y)$ at the witness $0$ without the cut of
\cref{claus:lem:struct}; \cref{comp:sec:clausal} gives the same derivation
for \pkg{uniform\_step}. Let $q\equiv\It{\Suc 0}{0}{u}{u}$, which is not
literally $0$ and satisfies $q=0$ by \cref{claus:lem:c0}(ii). Given
$\Gamma,E(0,0)\proves E(0,\Suc z)$, apply \rulename{Close} to the matrix
$E(0,y)$ with an unused first variable at the witness $q$; this is case
mode, its premise is the given derivation, and it gives
$\D{v}{y}E(0,y)$. \rulename{Open} adds the clause
$[E(0,0);\,y.\,E(0,y)]$. Now apply \rulename{Close} to $E(x,y)$ at the
witness $q$: its assumption $E(q,0)$ transports to $E(0,0)$,
\rulename{Use} gives $E(0,\Suc z)$, and transport gives $E(q,\Suc z)$.
\end{remark}

\section{Variants of the introduction}
\label{claus:sec:intro}

This section draws on \pkg{clause\_guards} and \pkg{rootfree\_clausal}.
The introduction of $\mathsf{C}$ selects its mode by the root of the matrix
and by the literal witness. A substitution can change the second test
(\cref{claus:prop:subst}). This section gives introductions whose mode is
selected by a guard on the matrix alone, and one whose mode is selected by
the instantiated zero base.

\subsection{The diagonal guard}
\label{claus:sec:diagonal}

This subsection draws on \pkg{clause\_guards}. The diagonal shape, its
instance at zero, its lift and its stability under substitution are those
of the diagonal introduction with native contexts
(\cref{guard:def:shapes,guard:def:selector,guard:thm:diagonal}); what
follows places that introduction in the clause calculus $\mathsf{C}$.

\begin{definition}[The calculus $\mathsf{K}$]
\label{claus:def:k}
Write $\operatorname{diag}_y(B)$ if the matrix $B(x,y)$ is diagonal in $y$
(\cref{guard:def:shapes}(c)), that is, literally
$\It{y}{y}{u}{h}=\It{y}{y}{v}{k}$ with $y$ the inner variable of the
enclosing $\D{x}{y}$ and $h,k$ arbitrary. The test is made on the
matrix before the witness is substituted. $\mathsf{K}$ is $\mathsf{C}$ with
the introduction
\[
 \infrule{\Gamma,B(t,\eta_B(z))\proves B(t,\Suc z)}
         {\Gamma\proves\D{x}{y}B}\ \rulename{Intro},
 \qquad
 \eta_B(z)\defeq
 \begin{cases}
 z,&\operatorname{diag}_y(B),\\
 0,&\text{otherwise,}
 \end{cases}
\]
with $z$ fresh for $\Gamma$, $t$ and $B$.
\end{definition}

No rule of $\mathsf{K}$ inspects a witness. \Cref{claus:lem:struct} and
\cref{claus:lem:eq} hold in
$\mathsf{K}$ with the same proofs: they use \rulename{Refl},
\rulename{Eval} and \rulename{Use} only, and injective renaming preserves
the diagonal test.

\begin{lemma}[Derived rules of $\mathsf{K}$]
\label{claus:lem:k}
\begin{enumerate}[label=(\roman*)]
\item If $\operatorname{diag}_y(B)$ then $\Gamma\proves B(t,0)$ for every
$t$ and $\Gamma$.
\item The introduction by cases of \cref{claus:lem:cases} is admissible for
every matrix and every witness.
\item The closure interfaces (i) to (iii) of \cref{claus:lem:closure},
numerical cases for every formula, and transport in every formula
(\cref{claus:lem:ftransport}) hold.
\end{enumerate}
\end{lemma}
\begin{proof}
(i) This is condition (S1) of \cref{guard:thm:diagonal}(i): the instance
is $\It{0}{0}{u}{h'}=\It{0}{0}{v}{k'}$, both sides equal $0$ by the zero
computation, and symmetry and transitivity join them. These are the
derivations of \cref{claus:lem:eq}, which read no clause.

(ii) For a matrix that is not diagonal this is \rulename{Intro}. For a
diagonal matrix cut the theorem $B(t,0)$ of (i) into the given branch
$\Gamma,B(t,0)\proves B(t,\Suc z)$. The result
$\Gamma\proves B(t,\Suc z)$ is weakened by the assumption $B(t,z)$, and
\rulename{Intro} in progress mode applies.

(iii) The proofs of \cref{claus:lem:closure}(i) to (iii) and of
\cref{claus:lem:ftransport} use the introduction by cases, \rulename{Use}
and \rulename{Elim}. Numerical cases follow from interfaces (i) and (ii).
\end{proof}

Equation induction is not immediate in $\mathsf{K}$, since progress is
restricted to diagonal matrices. For a literal equation
$E(y)\equiv(s(y)=t(y))$ let $H_E(y)\equiv\Delta_yE$ be the diagonal lift
of \cref{guard:def:shapes},
\[
 H_E(y)\defeq\bigl(L_s(y)=L_t(y)\bigr),\qquad
 L_s(y)\defeq\It{y}{y}{u}{s(y)},\quad L_t(y)\defeq\It{y}{y}{u}{t(y)},
\]
with $u$ not in $s,t$. The matrix $H_E$ is diagonal and $H_E(0)$ is a
theorem by \cref{claus:lem:k}(i). The successor computation gives
$L_s(\Suc z)=s(\Suc z)$ and $L_t(\Suc z)=t(\Suc z)$ (condition (S2) of
\cref{guard:thm:diagonal}(i)), hence
\begin{equation}
\label{claus:eq:liftsucc}
 H_E(\Suc z)\dashv\proves E(\Suc z).
\end{equation}
Numerical cases prove the two equations of the equational decoder that
follows \cref{guard:thm:diagonal},
\begin{equation}
\label{claus:eq:selrecover}
 \sel(y,s(0),L_s(y))=s(y),\qquad \sel(y,t(0),L_t(y))=t(y)
\end{equation}
without induction: at $0$ both sides of the first are $s(0)$ after the
zero computation, and at $\Suc z$ the selector and the iterator with
constant body compute to $s(\Suc z)$; likewise for $t$. Under $E(0)$,
congruence along $s(0)=t(0)$ and along $H_E(y)$ joins the left sides of
\eqref{claus:eq:selrecover}, so
\begin{equation}
\label{claus:eq:recovery}
 E(0),\ H_E(y)\proves E(y).
\end{equation}

\begin{theorem}[Equation induction in $\mathsf{K}$]
\label{claus:thm:keqind}
For every literal equation $E(y)$, every context $\Gamma$ and every term
$q$, derivations of $\Gamma\proves E(0)$ and of
$\Gamma,E(z)\proves E(\Suc z)$, $z$ fresh, yield a derivation of
$\Gamma\proves E(q)$. Consequently interface (iv) of
\cref{claus:lem:closure} and \cref{claus:cor:induction} hold in
$\mathsf{K}$.
\end{theorem}
\begin{proof}
Under the assumption $H_E(z)$, \eqref{claus:eq:recovery} and the given
derivation of $E(0)$ give $E(z)$. Cut this into the given step derivation
at its own variable $z$; the result is $E(\Suc z)$, and
\eqref{claus:eq:liftsucc} gives $H_E(\Suc z)$. \rulename{Intro} on the
diagonal matrix $H_E$ with an unused first variable proves $q_yH_E$. Its
zero instance is a theorem, so interface (ii) gives $H_E(q)$, and
\eqref{claus:eq:recovery} gives $E(q)$. No term is substituted in the
given step derivation, and the induction that is being derived is not used
for \eqref{claus:eq:selrecover}.

For interface (iv), work in $\Gamma,E(0)$. The first part gives
$E(\Suc w)$ for a fresh $w$, and interface (i) gives
$\Gamma\proves q_yE$. With a witness, \cref{claus:lem:k}(ii) gives the
third statement of \cref{claus:cor:induction}.
\end{proof}

\begin{theorem}[The diagonal clause calculus and $\Nat$]
\label{claus:thm:k}
The four statements of \cref{claus:thm:five} hold for $\mathsf{K}$.
\end{theorem}
\begin{proof}
(i) is proved as there, except for \rulename{Intro}. An introduction with
a matrix that is not diagonal is the introduction by cases of $\Nat$. A
diagonal matrix is a literal equation $E(x,y)$; the translated premise is
$\Gamma^{\natural},E(t,z)\proves E(t,\Suc z)$, the introduction of $Q$ by
progress gives $\Q{y}E(t,y)$, and the introduction by cases at $t$ with \rulename{QE}
gives $\D{x}{y}E$. (ii) is proved as there, with \cref{claus:lem:k}(ii) for
the introduction by cases and \cref{claus:thm:keqind} for the introduction
of $Q$ by progress. (iii) and (iv) follow as there.
\end{proof}

If $\sigma$ substitutes terms for free variables of a matrix $B(x,y)$,
lifted under $x$ and $y$, then $\operatorname{diag}_y(B\sigma)$ if and only
if $\operatorname{diag}_y(B)$: the root is preserved and reflected, and so
is the diagonal shape of a literal equation
(\cref{guard:prop:subst-shapes}(i) and (iii)).

\begin{theorem}[Substitution preserves the derivation tree]
\label{claus:thm:ksubst}
If $\Gamma\proves_{\mathsf{K}}A$ and $\sigma$ is a substitution of terms
for free variables, then $\Gamma\sigma\proves_{\mathsf{K}}A\sigma$ by a
derivation with the same tree of rules, the same selected clause
occurrences and the same evaluation rows, after a consistent renaming of
eigenvariables.
\end{theorem}
\begin{proof}
Induction on the derivation. \rulename{Refl} is immediate. The literal
matches of each row of $\Phi_2$ are preserved and substitution commutes with
the instantiation of the body. \rulename{Use} keeps its occurrence and its
antecedent. \rulename{Elim} uses a fresh variable in the substituted
context. For \rulename{Intro}, \cref{guard:prop:subst-shapes}(i) and (iii)
fix the mode, so the substituted branch is the branch that the same rule
requires. The argument is that of \cref{guard:thm:diagonal}(iii) and
\cref{guard:prop:np-subst}(i); the clause setting adds only that
\rulename{Use} and \rulename{Elim} commute with $\sigma$ and keep the
selected occurrence.
\end{proof}

\begin{remark}
\label{claus:rem:kpkgs}
In the Lean development $\mathsf{K}$ is \Lean{KPrf}, the equivalences are
\Lean{kprf\_iff\_nprf} and \Lean{kprf\_iff\_nprf\_of\_qfree}, and
\cref{claus:thm:ksubst} is \Lean{KPrf.subst}
(\Lean{NativeClausalDiagonal}). \pkg{clause\_guards} compares $\mathsf{K}$
with the six-rule native calculus that has the diagonal introduction
(\Lean{SPrf}), for which the translation $\tau$ preserves and reflects the
diagonal shape, so that introductions are copied; the seven-rule native
calculus with the same introduction is $\Nat_{\mathrm{diag}}$ of
\cref{guard:thm:diagonal}. $\mathsf{C}$ has
substitution as an admissible rule (\cref{claus:prop:subst}); the
statement for $\mathsf{K}$ is that the derivation is not repaired.
\Cref{claus:thm:keqind} differs from the derivation of \rulename{QI} in
\cref{guard:thm:selector}: it uses the decoder \eqref{claus:eq:recovery}
and the closure interfaces of the clause calculus in place of the
implication between equations of \cref{guard:lem:np-imp}.
\end{remark}

\subsection{Guards with a derivable base}
\label{claus:sec:guards}

This subsection draws on \pkg{clause\_guards}; the method of repairing an
introduction after a one-way change of guard is that of
\pkg{reflexive\_base}. It is the clause form of the selector criterion
with native contexts (\cref{guard:def:selector,guard:thm:selector,guard:prop:np-subst}):
conditions (a) and (b) below are the requirement on a selector and (S1),
and (c) is (S2) together with a decoder of the lift.

Let $G_y(B)$ be a decidable predicate on matrices that is invariant under
renaming of bound variables and under injective renaming of free
variables. $\mathsf{K}_G$ is $\mathsf{K}$ with $G$ in the place of
$\operatorname{diag}$. Consider the following conditions.
\begin{enumerate}[label=(\alph*)]
\item Every matrix that satisfies $G$ is a literal equation.
\item For every matrix $B$ that satisfies $G$ and every $t$, the zero
instance $B(t,0)$ is derivable by \rulename{Refl} and \rulename{Eval}.
\item For every literal equation $E(y)$ there is a matrix $H_E(y)$ that
satisfies $G$, with derivations of $H_E(\Suc z)\dashv\proves E(\Suc z)$ and
of $E(0),H_E(y)\proves E(y)$ that use \rulename{Refl}, \rulename{Eval},
\rulename{Use}, \rulename{Elim} and introductions by cases only.
\end{enumerate}

\begin{theorem}[Guard criterion]
\label{claus:thm:guard}
Under (a), (b) and (c) the four statements of \cref{claus:thm:five} hold
for $\mathsf{K}_G$. If moreover $G$ is preserved by substitution of terms
for free variables of the matrix, substitution is admissible in
$\mathsf{K}_G$ by a direct transformation of derivations that inserts
nodes only at introductions whose matrix acquires the guard. If $G$ is
also reflected by substitution, the transformation preserves the
derivation tree.
\end{theorem}
\begin{proof}
Condition (b) gives the introduction by cases for guarded matrices as in
\cref{claus:lem:k}(ii), hence the closure interfaces, numerical cases and
transport in formulas. Condition (c) and the proof of
\cref{claus:thm:keqind}, with the given lift in the place of the diagonal
lift, give equation induction and interface (iv). A guarded introduction
is sound in $\Nat$ because by (a) it is an instance of progress for a
literal equation; an unguarded one is an introduction by cases. The proof
of \cref{claus:thm:k} then applies.

For substitution induct on the derivation. A guarded matrix remains
guarded, and an unguarded matrix that remains unguarded keeps its rule.
If an unguarded matrix becomes guarded, the substituted branch assumes
the substituted zero instance; (b) proves that instance, cut removes the
assumption, the branch is weakened by the progress assumption, and the
guarded \rulename{Intro} applies. If $G$ is reflected this case does not
occur.
\end{proof}

\begin{proposition}[The mirror guard on the full language]
\label{claus:prop:mirror}
Let $\operatorname{mirror}_y(B)$ hold if $B\equiv(a=\It{y}{a}{u}{h})$.
The guard satisfies (a), (b) and (c) and is preserved by substitution. It
is not reflected by substitution.
\end{proposition}
\begin{proof}
The guard is the mirrored shape of \cref{guard:def:shapes}(b). (a) holds
by definition. (b) is condition (S1) of
\cref{guard:thm:mirror-selector}(i): the zero instance is
$a'=\It{0}{a'}{u}{h'}$, the symmetric form of the zero computation. (c):
for $E(y)\equiv(s(y)=t(y))$ the lift is the mirror gate
$M_E(y)\equiv H_yE\equiv\bigl(s(y)=\It{y}{s(y)}{u}{t(y)}\bigr)$,
$u\notin\FV(t)$, whose successor instance computes to $E(\Suc z)$, and
$E(0),M_E(y)\proves E(y)$ is the equational decoder of
\cref{guard:lem:decoder}. Both derivations use the computation
equations, equality and numerical cases only, which are available in the
clause calculus by \cref{claus:lem:eq} and interfaces (i) and (ii) of
\cref{claus:lem:closure}, the latter from introductions by cases,
\rulename{Use} and \rulename{Elim}. Preservation and the failure of
reflection are \cref{guard:prop:subst-shapes}(ii); for distinct free
variables $p,q$ the matrix $p=\It{y}{q}{u}{h}$ is not guarded, and the
substitution of $p$ for $q$ makes it guarded.
\end{proof}

Every matrix that is not guarded takes the case branch, so the mirror
calculus is defined on the whole $Q$-free language, as the mirror selector
with native contexts is on the full native language
(\cref{guard:thm:mirror-selector}); \pkg{mirror\_guard} restricts the
public language instead (\cref{guard:thm:mirror}).
% ed.: the earlier text attributed the public-language restriction to
% mirror_core; per guard:sec:full and guard:sec:mirror it is mirror_guard.
By \cref{claus:thm:guard} the mirror
calculus is equivalent to $\Nat$ and has admissible substitution with a
local repair, and it lacks the exact statement of
\cref{claus:thm:ksubst}. A guard that tests one iterator side without its
relation to the initializer does not satisfy (b).

\subsection{The fixed zero base on zero-normal terms}
\label{claus:sec:fixedzero}

This subsection draws on \pkg{rootfree\_clausal}; the guard is that of
\pkg{reflexive\_base} (\cref{comp:def:reflexive}), restricted from an
arbitrary literally reflexive base $r=r$ to the fixed base $0=0$, and
placed in the clause calculus $\mathsf{C}^{0}$.

\begin{definition}[The calculus $\mathsf{C}^{0}_{\mathrm{b}}$]
\label{claus:def:c0b}
$\mathsf{C}^{0}_{\mathrm{b}}$ is $\mathsf{C}^{0}$ with the introduction
\[
 \infrule{\Gamma,B\langle t,\eta\rangle\proves B\langle t,\Suc z\rangle}
         {\Gamma\proves\D{x}{y}B}\ \rulename{Intro}
\]
in which
\[
 \eta\defeq
 \begin{cases}
 z,&B\langle t,0\rangle\equiv(0=0),\\
 0,&\text{otherwise.}
 \end{cases}
\]
\end{definition}

The test is one comparison of the computed base with a fixed formula.
Substitution and normalization preserve the root of a formula, so the test
excludes matrices with root $D$; there is no separate test of the root or
of the witness.

\begin{theorem}[The fixed zero base]
\label{claus:thm:fixedzero}
For every zero-normal clause context $\Gamma$ and formula $F$,
$\Gamma\proves_{\mathsf{C}^{0}_{\mathrm{b}}}F$ if and only if
$\Gamma\proves_{\mathsf{C}^{0}}F$. Hence \cref{claus:thm:zero}(iii) and
(iv) hold for $\mathsf{C}^{0}_{\mathrm{b}}$.
\end{theorem}
\begin{proof}
The two calculi differ in the introduction only. \Cref{claus:lem:c0}(i) to
(v) hold in $\mathsf{C}^{0}_{\mathrm{b}}$ with the same proofs.

\emph{Cases in $\mathsf{C}^{0}_{\mathrm{b}}$.} Let
$\Gamma,B\langle t,0\rangle\proves B\langle t,\Suc z\rangle$. If the base
is not $0=0$ the rule applies. If it is $0=0$, cut \rulename{Refl} into
the branch, weaken by the assumption $B\langle t,z\rangle$ and apply the
rule in progress mode, as in \cref{comp:lem:rcases}. No auxiliary witness
is needed. The closure
interfaces (i) to (iii) and transport in formulas follow.

\emph{Equation progress in $\mathsf{C}^{0}_{\mathrm{b}}$.} Let
$E(y)\equiv(s(y)=t(y))$ and let $\Gamma,E(z)\proves E(\Suc z)$ be given;
the conclusion is $\Gamma\proves q_yE$. The derivation is that of
\cref{comp:thm:rprogress}, with the progressive lift
$H(y)\equiv\bigl(\sel(y,0,s(y))=\sel(y,0,t(y))\bigr)$ and the two
decoders $\mathrm{Dec}_{s,s_0}$ and $\mathrm{Dec}_{t,s_0}$ of
\cref{comp:lem:decoder}, $s_0\equiv s\langle 0\rangle$, the selector on
zero-normal terms being $\sel(v,a,b)\equiv\overline{I}(v,a;\,u.\,b)$. It
transfers to the clause calculus as follows. The closure by cases
\eqref{comp:eq:rqcase} is interface (i) of \cref{claus:lem:closure} with
$q_y$ in the place of $\Q{y}$, by the introduction by cases just derived;
\rulename{QE} is interface (ii); equality, congruence, weakening and cut
are \cref{claus:lem:c0}(i) to (v). The zero instance $H\langle 0\rangle$
is literally $0=0$, so the introduction of $\mathsf{C}^{0}_{\mathrm{b}}$ on
the matrix $H$ with an unused first variable is in progress mode, as
\rulename{RI} is there, and it concludes $q_yH$ directly, without the
\rulename{DE} step of that proof. As there, nothing is substituted into
the given derivation and no equation induction is used before the guarded
rule is applied.

\emph{From $\mathsf{C}^{0}$ to $\mathsf{C}^{0}_{\mathrm{b}}$.} An
\rulename{Intro} of $\mathsf{C}^{0}$ in case mode is an introduction by
cases. In progress mode the matrix is a literal equation and the witness
is $0$; equation progress gives $q_yB\langle 0,y\rangle$, and the
introduction by cases at the witness $0$ with interface (ii) gives
$\D{x}{y}B$.

\emph{From $\mathsf{C}^{0}_{\mathrm{b}}$ to $\mathsf{C}^{0}$.} If the
guard is false the rule is an introduction by cases
(\cref{claus:lem:c0}(vi)). If it is true the matrix is a literal equation,
interface (iv) of $\mathsf{C}^{0}$ applied to the matrix
$B\langle t,y\rangle$ gives $q_yB\langle t,y\rangle$, its base is $0=0$,
and interface (ii) gives $B\langle t,\Suc z\rangle$; the introduction by
cases at the witness $t$ concludes $\D{x}{y}B$.
\end{proof}

If the base is literally $0=0$ it remains so under substitution, since
$\nf_0(\nf_0(G)\sigma)\equiv\nf_0(G\sigma)$. A base that is not $0=0$ can become
$0=0$, which changes an introduction by cases into a progressive one; the
derived introduction by cases repairs that instance, as in
\cref{comp:lem:rsubst}. So the guard is preserved and not reflected, as
for the mirror guard.

\section{Variants of the context}
\label{claus:sec:context}

This section draws on \pkg{prefix\_clauses}, \pkg{clause\_contexts} and
\pkg{one\_body}. A context entry of $\mathsf{C}$ has an antecedent and a body. In
$\mathsf{V}$ the antecedent is optional (\cref{claus:sec:uniform}). This
section gives the representation of a formula by its prefix, with clauses
that keep one antecedent, and two calculi whose context entries have no
antecedent field at all.

\subsection{The prefix representation}
\label{claus:sec:prefix}

This subsection draws on \pkg{prefix\_clauses} and on the ranked
representation of \pkg{clause\_contexts}.

\begin{lemma}[Unique prefix]
\label{claus:lem:prefix}
Every $Q$-free formula is, in exactly one way up to the names of bound
variables,
\[
 \D{x_1}{y_1}\cdots\D{x_k}{y_k}(s=t),\qquad k\ge 0 .
\]
The number $k$ is unchanged by substitution of terms. A native formula has
this form if and only if it is $Q$-free.
\end{lemma}
\begin{proof}
An equation has $k=0$. A formula $\D{x}{y}B$ has exactly one formula child;
the grammar has no branching constructor, so removing binders reaches
exactly one equation, and restoring them in order reconstructs the
formula. Substitution acts on the two terms of the leaf.
\end{proof}

The formula is therefore the record $[k,s,t]$, written $D^k(s=t)$: a
number and two terms, with $2k$ bound variables in scope at the leaf. Only
term syntax is recursive. The record is a bijective encoding of the
$Q$-free formulas; it is not a prenex normal form theorem for ordinary
formulas, and each $D$ keeps its conditional reading. With de Bruijn
indices starting at $1$, the two outer variables of the matrix of
$D^{k+1}(s=t)$ are the indices $2k+2$ (the witness) and $2k+1$ at the leaf,
and an iterator body adds one to the index sought.

\begin{lemma}[Opening commutes with packing]
\label{claus:lem:packing}
The record of an instance $B(t,r)$ of the matrix of $D^{k+1}(s=t)$ is
obtained from $[k,s,t]$ by substituting at the index $2k+2$ and then at the
index $2k+1$, with larger indices decreased when a variable is removed and
the index sought increased by one under each iterator body.
\end{lemma}
\begin{proof}
Each removed pair of binders adds two to the depth at which the
substitution acts. At the leaf induct on the two terms: a variable is
replaced at the indicated index, $\Suc$ adds no binder and the body of an
iterator adds one. These are the clauses of capture-avoiding substitution
on the nested representation.
\end{proof}

In the Lean development the predicate \Lean{NFm.Pre k A} states that $A$ is
$k$ binders $D$ followed by one equation, with \Lean{NFm.qfree\_iff\_pre},
\Lean{NFm.Pre.unique} and \Lean{NFm.Pre.subst}; \Lean{VPrf.pre} states that
$\mathsf{V}$ derives only such formulas from $Q$-free assumptions
(\Lean{NativePrefix}).

\begin{definition}[The calculus $\mathsf{PF}$]
\label{claus:def:pf}
Formulas are records $D^k(s=t)$ on raw terms. A formula is
\emph{reflexive at its leaf} if it is $D^k(r=r)$ with literally identical
terms. Clauses are $[A;\,y.\,B(y)]$ with one antecedent, an assumption $A$
is $[0=0;\,v.\,A]$, and every context contains the clause
\[
 \rho\defeq[\,0=0;\,y.\,y=y\,].
\]
The rules are \rulename{Use}, \rulename{Eval} with $\Phi_2$,
\rulename{Elim} and \rulename{Intro} of $\mathsf{C}$, with one change: in
\rulename{Use} of $[A;\,y.\,B]$ the premise is omitted when $A$ is
reflexive at its leaf. In \rulename{Intro} the test that the matrix is a
literal equation is the test that the conclusion has $k=1$.
\end{definition}

There are four rule families and the fixed clause $\rho$. The test of
reflexivity at the leaf compares two term trees and evaluates nothing.

\begin{lemma}[Reflexive prefixes]
\label{claus:lem:reflprefix}
Every formula $D^k(r=r)$ is derivable in every context in $\mathsf{C}$, in
$\mathsf{V}$ with $\mathcal{R}$, and in $\mathsf{PF}$.
\end{lemma}
\begin{proof}
Induction on $k$. For $k=0$ use \rulename{Refl} in $\mathsf{C}$,
\rulename{Use} of $\mathcal{R}$ in $\mathsf{V}$, and \rulename{Use} of
$\rho$ in $\mathsf{PF}$, where the antecedent $0=0$ is reflexive and the
premise is omitted. For $k>0$ apply \rulename{Intro} at the witness
$\Suc 0$, which is case mode whatever the matrix. The target at the fresh
successor has $k-1$ binders and identical terms at its leaf, since both
receive the same substitutions; the induction hypothesis proves it, and
the added clause is not used.
\end{proof}

\begin{theorem}[The prefix calculus]
\label{claus:thm:prefix}
For every clause context $\Gamma$ and formula $A$, with formulas packed on
the left,
\[
 \rho,\Gamma\proves_{\mathsf{PF}}A\quad\Longleftrightarrow\quad
 \Gamma\proves_{\mathsf{C}}A .
\]
Hence $\rho,\operatorname{ord}(\Delta)\proves_{\mathsf{PF}}F$ if and only
if $\Delta\proves_{\Nat}F$, for $Q$-free $\Delta$ and $F$. Both formula
round trips between records and nested formulas are identities.
\end{theorem}
\begin{proof}
From right to left: pack all formulas and add $\rho$. \rulename{Refl} at
$t$ becomes \rulename{Use} of $\rho$ at $t$ without premise. The other
rules translate directly; by \cref{claus:lem:packing} substitutions agree,
and the free variables are unchanged, so the eigenvariable conditions
agree. From left to right: a \rulename{Use} of the distinguished $\rho$
becomes \rulename{Refl}. Any other \rulename{Use} keeps its translated
premise; if the premise was omitted, \cref{claus:lem:reflprefix} supplies
a derivation of the antecedent in $\mathsf{C}$. \rulename{Eval},
\rulename{Elim} and \rulename{Intro} translate directly. No cut is needed.
The second statement is \cref{claus:thm:five}(iii).
\end{proof}

By \cref{claus:thm:uniform} the same holds against $\mathsf{V}$, with the
empty antecedent represented by $0=0$. \pkg{prefix\_clauses} states the
theorem against \Lean{CPrf} and \Lean{VPrf}.

\subsection{Rank from variable support}
\label{claus:sec:support}

This subsection draws on the optional notes of \pkg{clause\_contexts}. It
shows that the number $k$ of a public record can be recovered from the
variables of its leaf, at the price of identifying some theorems.

Work in $\mathsf{V}^{0}$ with records and de Bruijn indices. For a term $t$
let $\mu(t)$ be the greatest positive index that is free relative to the
iterator binders of $t$, and $0$ if there is none:
$\mu(j)=\max(j,0)$, $\mu(\Suc t)=\mu(t)$ and
$\mu(\It{n}{a}{u}{h})=\max(\mu(n),\mu(a),\max(\mu(h)-1,0))$. For a record
$F=[k,s,t]$ let $r(F)\defeq\lceil\max(\mu(s),\mu(t))/2\rceil$, so that
$r(F)\le k$. Define
\[
 P[k,s,t]\defeq
 \begin{cases}
 (s,t),&r(F)=k,\\
 (0,0),&r(F)<k,
 \end{cases}
 \qquad
 U(s,t)\defeq\bigl[\lceil\max(\mu(s),\mu(t))/2\rceil,\,s,\,t\bigr].
\]
The pair $(0,0)$ decodes to $0=0$. The quotient calculus has pairs as the
formulas of its judgments and records in its clauses, which keep their
rank: \rulename{Use} asks for derivations of $P(A_i)$ and returns
$P(B\langle t/y\rangle)$; \rulename{Step} acts on pairs of rank zero;
\rulename{Open} decodes a major premise of nonzero rank, adds its clause as
a record and derives the same pair in the branch; \rulename{Close} decodes
a goal of nonzero rank, adds its assumption as a record and asks for the
packed successor instance. $P$ is applied only to formulas all of whose
clause variables are instantiated, and it is recomputed after every
normal substitution.

\begin{proposition}[Support quotient]
\label{claus:prop:support}
For every record $F$, $F$ and $U(P(F))$ are interderivable in
$\mathsf{V}^{0}$ in every context that contains $\mathcal{R}$, and
$P(U(a))=a$ for every pair $a$. The quotient calculus derives $P(F)$ from
$\mathcal{R}$ if and only if $\mathsf{V}^{0}$ derives $F$ from
$\mathcal{R}$.
\end{proposition}
\begin{proof}
If $r(F)=k$ then $U(P(F))$ is $F$. If $r(F)<k$, the two variables of the
outermost binder, with the indices $2k$ and $2k-1$, do not occur in the
leaf, so $F\equiv\D{x}{y}B$ with $x,y$ not in $B$. \rulename{Close} at the
witness $0$ has $B$ both as its assumption and as its target in either
mode, so $F$ is a theorem; $0=0$ is a theorem; and each follows from the
other by weakening. $P(U(a))=a$ because $U$ uses the rank that $P$ tests.

Soundness of the quotient rules: by the first part a derivation of
$U(P(A))$ converts to one of $A$ and conversely, by cut and weakening.
\rulename{Use} converts its packed premises, applies \rulename{Use} of
$\mathsf{V}^{0}$ and converts the result. \rulename{Step} involves rank
zero only. \rulename{Close} converts the local base and the target and
applies \rulename{Close}; the decoded goal has the rank derived from its
support and packs to the same pair. \rulename{Open} decodes a record with
root $D$ and adds the clause of \rulename{Open}.

Completeness: transform a derivation of $\mathsf{V}^{0}$, packing every
conclusion. \rulename{Use} and \rulename{Step} transform directly. A
\rulename{Close} whose goal keeps its rank is the same \rulename{Close};
if the goal collapses, \rulename{Use} of $\mathcal{R}$ at $0$ proves
$0=0$ and the premise is discarded. \rulename{Open} transforms directly
when its major premise keeps its rank. If the major premise $\D{x}{y}B$
collapses, the added clause is $[B;\,y.\,B]$, and each use of it returns
the formula of its own premise; delete those uses and keep the premise
derivations, removing the clause from the branch. The branch may still
contain the eigenvariable $b$ in intermediate terms. To discharge it, use
the theorem $T\equiv\D{x}{y}(x=x)$, which has rank $1$ and support index
$2$ and does not collapse. \rulename{Close} at the witness $0$ proves $T$,
both instances being $0=0$. \rulename{Open} of $T$ at the fresh $b$ adds
the clause $[b=b;\,y.\,b=b]$, the transformed branch is weakened by it,
and the rule concludes the goal in the outer scope.
\end{proof}

The quotient contracts a binder only when both its variables are absent;
theorems such as $\D{x}{y}(x=x)$ are kept, and no completeness of the test
for theorems is claimed. A clause body has one more bound variable, the
clause variable, so the same support can belong to a body of rank zero
with a clause variable or to a formula of rank one; clause bodies therefore
keep an explicit rank.

\subsection{One body per context entry}
\label{claus:sec:onebody}

This subsection draws on \pkg{one\_body}. An opened occurrence
$[B(b,0);\,y.\,B(b,y)]$ is determined by its body. An assumption can be
stored in the same form.

\begin{definition}[Gate and the calculus $\mathsf{OB}$]
\label{claus:def:ob}
For $A\equiv D^d(s=t)$ and a variable $v$ outside the $2d$ bound variables
the \emph{gate} is
\[
 \operatorname{gate}_vA\defeq
 D^d\bigl(\sel(v,0,s)=\sel(v,0,t)\bigr).
\]
A context of $\mathsf{OB}$ is a finite list of \emph{bodies} $B(v)$, each
a formula with one distinguished variable; a body reads
$B(0)\to\forall v\,B(v)$, which is the reading of $\Q{v}B$. The rules are
\rulename{Refl} for equations, $\rulename{Eval}_0$ and
$\rulename{Eval}_{\Suc}$, and
\[
 \infrule{\Gamma\proves B(0)}{\Gamma\proves B(t)}\ \rulename{Use}
 \qquad
 \infrule{\Gamma\proves\D{x}{y}B(x,y)\quad\Gamma,B(b,y)\proves A}
         {\Gamma\proves A}\ \rulename{Open}
\]
\[
 \infrule{\Gamma,\operatorname{gate}_v\bigl(B(w,\eta)\bigr)\proves
          B(w,\Suc z)}
         {\Gamma\proves\D{x}{y}B(x,y)}\ \rulename{Intro}
\]
where \rulename{Use} selects a body $B(v)$ of $\Gamma$, \rulename{Open}
adds the body $y\mapsto B(b,y)$ with $b$ fresh, and in \rulename{Intro}
the variables $z$ and $v$ are fresh, $w$ is the witness, and $\eta$ is $z$
if $B$ is a literal equation and $w\equiv 0$, and $0$ otherwise.
\end{definition}

A context entry has no antecedent field, and there is no entry for an
ordinary assumption. The calculus is formalized in Lean as \Lean{OBPrf}
(\Lean{NativeOneBody}); its theorems are listed in
\cref{claus:rem:obpkgs}. The gate is a construction on the existing grammar.
Its variable $v$ lies outside the binders of $A$, and the terms of the
leaf are placed under the new iterator binder of the selector, which in a
de Bruijn representation requires one lifting.

The subsystem with \rulename{Refl} and the two evaluation rows gives the
equality reasoning of \cref{claus:lem:eq} in every context. Write $s\simeq t$ if this
subsystem proves $s=t$. Its derivations have no mode test and are stable
under substitution. In particular
\begin{equation}
\label{claus:eq:endpoints}
 \sel(0,0,s)\simeq 0,\qquad \sel(\Suc z,0,s)\simeq s .
\end{equation}
For formulas $A,A'$ with the same prefix, aligned, write $A\approx A'$ if
their left terms are related by $\simeq$ and their right terms are related
by $\simeq$, the bound variables of the prefix being treated as free. The
definition does not presuppose transport in formulas.

\begin{lemma}[Deep reflexivity]
\label{claus:lem:deeprefl}
Every $D^d(t=t)$ is derivable in $\mathsf{OB}$ in every context. Weakening
and injective renaming of free variables are admissible.
\end{lemma}
\begin{proof}
Induction on $d$ as in \cref{claus:lem:reflprefix}: at a binder apply
\rulename{Intro} at the witness $\Suc 0$ and ignore the new gate body.
Weakening and renaming are inductions on derivations.
\end{proof}

\begin{lemma}[Conversion and access]
\label{claus:lem:access}
For every $d$ the following hold in $\mathsf{OB}$.
\begin{description}
\item[$L_d$] If $A\approx A'$ and both have $d$ binders, a derivation of
either in a context converts to a derivation of the other in that context.
\item[$H_d$] If $A$ has $d$ binders and $\operatorname{gate}_vA$ is a body
of $\Gamma$, then $\Gamma\proves A$.
\end{description}
\end{lemma}
\begin{proof}
The statements are proved in the order $L_0,H_0,L_1,H_1,\dots$.

$L_0$ is symmetry and transitivity. For $L_d$ with $d>0$ assume $L_{d-1}$
and $H_{d-1}$, and let $A\equiv\D{x}{y}B$, $A'\equiv\D{x}{y}B'$. Apply
\rulename{Open} to the given derivation of $A$ at a fresh $b$, which adds
the body $B(b,y)$. Apply \rulename{Intro} for $A'$ at the witness $b$. A
variable is not the literal $0$, so this is case mode, and the branch has
the body $\operatorname{gate}_v B'(b,0)$. $H_{d-1}$ gives $B'(b,0)$.
$L_{d-1}$, with the endpoint derivations instantiated at $x\defeq b$,
$y\defeq 0$, gives $B(b,0)$. \rulename{Use} of the body $B(b,y)$ at
$\Suc z$ gives $B(b,\Suc z)$, and $L_{d-1}$ gives $B'(b,\Suc z)$.
\rulename{Intro} concludes $A'$ and \rulename{Open} discharges $b$. The
other direction uses the symmetric endpoint equations.

For $H_d$, given $L_d$: let $T_d\equiv D^d(0=0)$ with the prefix of $A$.
\Cref{claus:lem:deeprefl} proves $T_d$, and \eqref{claus:eq:endpoints}
gives $\operatorname{gate}_0A\approx T_d$ and
$\operatorname{gate}_1A\approx A$. $L_d$ gives
$\operatorname{gate}_0A$. \rulename{Use} of the body
$\operatorname{gate}_vA$ at $1$ gives $\operatorname{gate}_1A$, and $L_d$
gives $A$.
\end{proof}

Conversion at $d$ binders uses access at $d-1$ binders only, so the
induction is well founded. It is an induction on formula syntax; it adds
no induction to the arithmetic and does not presuppose the access that it
proves. The terms of the leaf may contain the bound variables of the
prefix, so the endpoint equations live under the binders and cannot be
replaced by one global substitution of equals.

\begin{theorem}[From clauses to bodies]
\label{claus:thm:ctob}
For $Q$-free $A_1,\dots,A_n,F$: if
$\operatorname{ord}(A_1),\dots,\operatorname{ord}(A_n)\proves_{\mathsf{C}}F$
then
$\operatorname{gate}A_1,\dots,\operatorname{gate}A_n\proves_{\mathsf{OB}}F$.
\end{theorem}
\begin{proof}
Induct on the derivation with its canonical annotation, translating an
assumption occurrence $\operatorname{ord}(A)$ to the body
$\operatorname{gate}_vA$ and an opened occurrence
$[B(b,0);\,y.\,B(b,y)]$ to the body $B(b,y)$. \rulename{Refl} and
\rulename{Eval} are unchanged. A \rulename{Use} of an assumption
occurrence is replaced by $H_d$ of \cref{claus:lem:access}; its premise
$0=0$ is discarded. A \rulename{Use} of an opened occurrence has the
premise $B(b,0)$ and is \rulename{Use} of $\mathsf{OB}$.
\rulename{Elim} adds an opened occurrence, as \rulename{Open} adds its
body. \rulename{Intro} adds the assumption $B(w,\eta)$, whose translation
is the gate body that \rulename{Intro} of $\mathsf{OB}$ adds; witness,
mode and eigenvariables are unchanged.
\end{proof}

The translation uses no cut and no weakening of $\mathsf{OB}$ beyond
\cref{claus:lem:deeprefl}, and no substitution in a derivation with
\rulename{Intro}.

\begin{lemma}[The gate in $\Nat$]
\label{claus:lem:gatenat}
For every $Q$-free $A$ and fresh $v$, in every context,
$A\proves_{\Nat}\Q{v}\operatorname{gate}_vA$ and
$\Q{v}\operatorname{gate}_vA\proves_{\Nat}A$.
\end{lemma}
\begin{proof}
The equations \eqref{claus:eq:endpoints} hold in $\Nat$, so
\cref{claus:lem:unq}(ii), applied along the prefix, gives
$\operatorname{gate}_{\Suc z}A\dashv\proves A$ and
$\operatorname{gate}_0A\dashv\proves D^d(0=0)$; the latter formula is
derivable by $d$ introductions by cases. From $A$ one obtains
$\operatorname{gate}_{\Suc z}A$, the introduction by cases with an unused
first variable gives $\D{w}{v}\operatorname{gate}_vA$, and
\cref{claus:lem:unq}(i) gives $\Q{v}\operatorname{gate}_vA$. Conversely \rulename{QE}
with the derivable base $\operatorname{gate}_0A$ gives
$\operatorname{gate}_1A$, hence $A$.
\end{proof}

\begin{theorem}[One-body contexts and $\Nat$]
\label{claus:thm:onebody}
If $\Gamma$ is a context of bodies and $\Gamma\proves_{\mathsf{OB}}A$, then
$\Gamma^{Q}\proves_{\Nat}A$, where $\Gamma^{Q}$ replaces each body $B(v)$
by $\Q{v}B$. For $Q$-free $A_1,\dots,A_n,F$ the following three statements
are equivalent:
\begin{align*}
 &\operatorname{gate}A_1,\dots,\operatorname{gate}A_n
   \proves_{\mathsf{OB}}F,\\
 &A_1,\dots,A_n\proves_{\Nat}F,\\
 &\operatorname{ord}(A_1),\dots,\operatorname{ord}(A_n)
   \proves_{\mathsf{C}}F .
\end{align*}
\end{theorem}
\begin{proof}
Soundness is by induction on the derivation. \rulename{Refl} and
\rulename{Eval} are as in \cref{claus:thm:five}(i). \rulename{Use} is \rulename{QE}
on the assumption $\Q{v}B$. \rulename{Open} is \rulename{DE}, which adds
$\Q{y}B(b,y)$. For \rulename{Intro} the translated premise is
$\Gamma^{Q},\Q{v}\operatorname{gate}_v(B(w,\eta))\proves B(w,\Suc z)$;
\cref{claus:lem:gatenat} and cut give
$\Gamma^{Q},B(w,\eta)\proves B(w,\Suc z)$, and the two modes are treated
as in \cref{claus:thm:five}(i). For the equivalences: soundness and
\cref{claus:lem:gatenat} with cut lead from the first statement to the
second; the second and the third are equivalent by
\cref{claus:thm:five}(iii); and \cref{claus:thm:ctob} leads from the third
to the first.
\end{proof}

\begin{remark}
\label{claus:rem:obpkgs}
In the Lean development the first
equivalence is \Lean{obprf\_iff\_nprf}, its form for arbitrary native
sequents through $\tau$ is \Lean{obprf\_unQ\_iff\_nprf}, and the
equivalence with $\mathsf{C}$ is \Lean{obprf\_iff\_cprf};
\cref{claus:lem:access} is \Lean{OBPrf.conv} and
\Lean{OBPrf.acc}, and \cref{claus:thm:ctob} is \Lean{CPrf.toOB}. The Lean
proof derives the gate at $0$ directly (\Lean{OBPrf.gtZero}) and does not
convert it from deep reflexivity. \pkg{one\_body} proves soundness through
$\Teq$ and the bridge (\cref{base:thm:bridge}); the proof above reads a
body as a formula $\Q{v}B$ of $\Nat$, as the Lean proof does.
\end{remark}

\begin{remark}[A padding trap]
\label{claus:rem:padding}
To align the numbers of binders of two formulas, a formula $A$ is padded
as $\D{x}{y}\operatorname{gate}_yA$ with $x$ unused, which is equivalent
to $A$ by \cref{claus:lem:gatenat}. Padding by a binder both of whose
variables are unused is incorrect: $\D{x}{y}A$ then reads
$\exists x\,(A\to A)$ and is a theorem
(compare \cref{claus:prop:support}).
\end{remark}

\subsection{Homogeneous contexts with an asymmetric gate}
\label{claus:sec:homog}

This subsection draws on the appendix of \pkg{prefix\_clauses}. It removes
the antecedent field as $\mathsf{OB}$ does, keeps reflexivity as a fixed
entry, and uses a gate whose zero instance is literally reflexive.

\begin{definition}[The calculus $\mathsf{HB}$]
\label{claus:def:hb}
For a fresh variable $v$ let
\[
 \kappa_v(s=t)\defeq\bigl(\It{0}{s}{u}{t}=\It{v}{s}{u}{t}\bigr),
 \qquad
 \kappa_v(\D{x}{y}B)\defeq\D{x}{y}\kappa_vB,
\]
with $u$ not in $t$. A context is a list of bodies and always contains
the body $\rho(y)\equiv(y=y)$. The rules are \rulename{Eval} with $\Phi_2$,
and
\[
 \infrule{\Gamma\proves B(0)}{\Gamma\proves B(t)}
 \qquad
 \infrule{B(0)\equiv D^j(r=r)}{\Gamma\proves B(t)}
 \qquad\rulename{Use}
\]
for a body $B(y)$ of $\Gamma$, the rule \rulename{Open} of $\mathsf{OB}$,
and \rulename{Intro} of $\mathsf{OB}$ with
$\kappa_v\bigl(B(t,\eta)\bigr)$ as the added body.
\end{definition}

There are four rule families, no rule \rulename{Refl} and no hypothesis
rule. The second form of \rulename{Use} omits the premise exactly when the
raw base $B(0)$ is reflexive at its leaf, for some $j\ge 0$; a theorem
with root $D$ of another shape does not qualify. The gate satisfies
\[
 \kappa_0(s=t)\ \text{is literally reflexive},\qquad
 \kappa_{\Suc z}(s=t)\dashv\proves(s=t),
\]
the second by the two computation equations, since $u$ is not in $t$. The
gate changes neither the number of binders nor the root of a formula.

\begin{lemma}[Gate retrieval]
\label{claus:lem:retrieval}
In $\mathsf{HB}$, in every context:
\begin{enumerate}[label=(\roman*)]
\item $t=t$ is derivable, every $D^j(r=r)$ is derivable, and the equality
rules of \cref{claus:lem:eq} hold;
\item $\kappa_0A$ is derivable for every $A$;
\item a derivation of $\kappa_1A$ converts to a derivation of $A$, and
conversely;
\item if $\kappa_vA$ is a body of $\Gamma$ then $\Gamma\proves A$.
\end{enumerate}
\end{lemma}
\begin{proof}
(i) \rulename{Use} of $\rho$ at $t$ has the base $0=0$, so its premise is
omitted and it concludes $t=t$. Reflexive prefixes are proved as in
\cref{claus:lem:reflprefix}, with \rulename{Intro} at the witness
$\Suc 0$ and the new gate body ignored. The equality rules use
reflexivity and \rulename{Eval} only. (ii) $\kappa_0A$ has the prefix of
$A$ and identical terms at its leaf, so it is a reflexive prefix.

(iii) and (iv) are proved together by induction on the number of binders
of $A$. For an equation, (iii) is the two computations and transport. For
(iv), \rulename{Use} of the body $\kappa_vA$ at $1$ has a base that is
reflexive at its leaf, so its premise is omitted and it concludes
$\kappa_1A$; (iii) at the same number of binders gives $A$. For (iii) at
$A\equiv\D{x}{y}B$, assume (iii) and (iv) for formulas with fewer binders.
Given $\kappa_1A\equiv\D{x}{y}\kappa_1B$, apply \rulename{Open} at a fresh
$b$, which adds the body $\kappa_1B(b,y)$, and \rulename{Intro} for $A$ at
the witness $b$, which is case mode and adds the body
$\kappa_v(B(b,0))$. (iv) gives $B(b,0)$, (iii) gives $\kappa_1B(b,0)$,
\rulename{Use} of the body $\kappa_1B(b,y)$ at $\Suc z$ gives
$\kappa_1B(b,\Suc z)$, and (iii) gives $B(b,\Suc z)$. The converse opens
$A$, introduces $\kappa_1A$ at the witness $b$, retrieves the local zero
assumption by (iv), converts it by (iii), uses the body $B(b,y)$ at
$\Suc z$ and converts the result by (iii).
\end{proof}

\begin{theorem}[Homogeneous contexts and $\Nat$]
\label{claus:thm:homog}
For $Q$-free $A_1,\dots,A_n,F$,
\[
 \rho,\kappa A_1,\dots,\kappa A_n\proves_{\mathsf{HB}}F
 \quad\Longleftrightarrow\quad
 A_1,\dots,A_n\proves_{\Nat}F .
\]
\end{theorem}
\begin{proof}
From right to left translate a derivation of $\mathsf{C}$
(\cref{claus:thm:five}(iii)) by the origins of its occurrences, as in
\cref{claus:thm:ctob}: \rulename{Refl} becomes \rulename{Use} of $\rho$, a
\rulename{Use} of an assumption occurrence becomes
\cref{claus:lem:retrieval}(iv), a \rulename{Use} of an opened occurrence
becomes \rulename{Use} with its premise, and \rulename{Elim} and
\rulename{Intro} become \rulename{Open} and \rulename{Intro}.

From left to right read a body $B(y)$ as $\Q{y}B$ and induct as in
\cref{claus:thm:onebody}. The reading $\Q{y}(y=y)$ of $\rho$ is derivable
in $\Nat$ by progress. A \rulename{Use} with omitted premise is \rulename{QE} whose
base $D^j(r=r)$ is derivable in $\Nat$ by $j$ introductions by cases. For
\rulename{Intro} use $A\proves_{\Nat}\Q{v}\kappa_vA$, which is proved as
\cref{claus:lem:gatenat} from $\kappa_{\Suc z}A\dashv\proves A$. Finally
$\Q{v}\kappa_vA_i$ follows from $A_i$, and cut gives the sequent on the
right.
\end{proof}

\begin{remark}
\label{claus:rem:hbpkgs}
\pkg{prefix\_clauses} proves completeness from the seven-rule witness
calculus of \pkg{witness\_rule} by the same strengthened induction over
contexts, and soundness through $\Teq$. The package states $\mathsf{HB}$
on records; the gate leaves $k$ unchanged and acts on the leaf. This
second theorem of the package is not formalized in Lean.
$\mathsf{OB}$ and $\mathsf{HB}$ differ in three points: the gate
(symmetric selector or asymmetric iterator), reflexivity (a rule for
equations or the entry $\rho$ with omitted premises), and the proof of
access (conversion along $\approx$ from deep reflexivity or a literally
reflexive zero gate).
\end{remark}

\section{Fusion: evaluation at the discharge}
\label{claus:sec:fusion}

This section draws on \pkg{evaluation\_discharge}, and
\cref{claus:sec:pd} on \pkg{476\_byte\_presentation}. The two rows of
\rulename{Eval} and the elimination of $D$ become one rule family, which
leaves three: \rulename{Use}, \rulename{Intro} and \rulename{Discharge}.

\subsection{The raw calculus}

\begin{definition}[One evaluation step]
\label{claus:def:rho}
For formulas $E$ and $F$ the formula $\varrho_E(F)$ is defined by three
disjoint cases, every match being literal:
\begin{align*}
 \varrho_{\,s=t}\bigl(\It{0}{s}{u}{h}=r\bigr)&\defeq(r=t),\\
 \varrho_{\,\It{c}{a}{u}{h}=t}\bigl(\It{\Suc c}{a}{u}{h}=r\bigr)
   &\defeq\bigl(r=h\subst{t}{u}\bigr),\\
 \varrho_E(F)&\defeq F\qquad\text{in every other case.}
\end{align*}
\end{definition}

The first case has a literal zero counter and the second a successor, so
they do not overlap, and $\varrho$ is a total function of $(E,F)$. The
replacement term, the body and the counter are read from the two
formulas. An $E$ with root $D$ always selects the third case. In the Lean
development the graph of $\varrho$ is a relation, with
\Lean{Rho.functional} and \Lean{Rho.total}.

\begin{definition}[The calculus $\mathsf{ED}$]
\label{claus:def:e3}
Terms are raw, contexts are lists of uniform clauses, and the rules are
\rulename{Use} and \rulename{Intro} of $\mathsf{V}$
(\cref{claus:def:v}) and
\[
 \infrule{\Gamma\proves\D{x}{y}B\quad
          \Gamma,\mathcal{B}_b\proves E\quad
          \Gamma,\mathcal{B}_b\proves F}
         {\Gamma\proves\varrho_E(F)}\ \rulename{Discharge}
\]
where $\mathcal{B}_b\defeq[B(b,0);\,y.\,B(b,y)]$ and $b$ is fresh for $\Gamma$ and $B$ and does not occur free in
$\varrho_E(F)$. The formulas $E$ and $F$ may contain $b$.
\end{definition}

The rule opens $\D{x}{y}B$ at $b$, proves $E$ and $F$ in that scope,
performs one evaluation step, and exports the result if the step has
removed $b$ or $b$ was absent. The calculus is formalized in Lean as
\Lean{E3Prf} (\Lean{NativeEvalDischarge}), with $\varrho$ as \Lean{Rho};
its theorems are listed in \cref{claus:rem:e3pkgs}. The third case does not derive $F$ without
a premise: the rule requires a derivation of $F$. The two patterns of
$\varrho$ and its identity case are part of the specification.

\begin{lemma}
\label{claus:lem:theta}
Weakening and injective renaming of free variables are admissible in
$\mathsf{ED}$ and in $\mathsf{V}$. The formula
$\Theta\equiv\D{x}{y}(0=0)$ is derivable in both, in every context, also in
one without $\mathcal{R}$.
\end{lemma}
\begin{proof}
Weakening and renaming are inductions on derivations with eigenvariables
made fresh; an injective renaming preserves and reflects the literal
matches of $\varrho$, the constructors $0$ and $\Suc$, and the mode of
\rulename{Intro}. For $\Theta$ apply \rulename{Intro} at the witness $0$.
The matrix is a literal equation, the assumption and the target are both
$0=0$, and \rulename{Use} of the added clause $\langle 0=0\rangle$ closes
the branch.
\end{proof}

\begin{theorem}[Evaluation at the discharge]
\label{claus:thm:fusion}
For every context $\Gamma$ of uniform clauses and every formula $A$,
\[
 \Gamma\proves_{\mathsf{ED}}A\quad\Longleftrightarrow\quad
 \Gamma\proves_{\mathsf{V}}A .
\]
The translations leave the end sequent unchanged, and $\mathcal{R}$ need
not be in $\Gamma$. Hence, for $Q$-free $\Delta$ and $F$,
$\mathcal{R},\langle\Delta\rangle\proves_{\mathsf{ED}}F$ if and only if
$\Delta\proves_{\Nat}F$.
\end{theorem}
\begin{proof}
From left to right translate the three premises of a
\rulename{Discharge}. If $\varrho_E(F)$ is given by the first or the
second case, apply the corresponding row of \rulename{Eval} to the
translated derivations of $E$ and $F$ in the context
$\Gamma,\mathcal{B}_b$. If it is given by the third case, the translated
derivation of $F$ has the required conclusion. Then apply \rulename{Elim};
its eigenvariable condition is the condition that $b$ is not free in
$\varrho_E(F)$. \rulename{Use} and \rulename{Intro} are unchanged.

From right to left translate \rulename{Elim} with major premise
$\D{x}{y}B$ and branch conclusion $A$ by \rulename{Discharge} with the
same major premise, with $E\defeq\D{x}{y}B$, derived in the branch by
weakening the translated major premise, and with $F\defeq A$. Since $E$
has root $D$, $\varrho_E(A)\equiv A$. Translate an \rulename{Eval} with
premises $E$ and $F$ by \rulename{Discharge} with the major premise
$\Theta$ of \cref{claus:lem:theta} and with the translated derivations of
$E$ and $F$ weakened to the context
$\Gamma,[0=0;\,y.\,0=0]$. The corresponding case of $\varrho$ matches and
gives the conclusion of the \rulename{Eval}, which contains no $b$. The
second statement is \cref{claus:thm:uniform}.
\end{proof}

Under the ordinary reading \rulename{Discharge} is sound directly:
eliminate the existential at $b$; in the first case the zero computation
and $s=t$ turn the second equation into $r=t$, in the second case the
successor computation and the equation for the preceding iterate give
$r=h\subst{t}{u}$, and the third case keeps the formula; the conclusion
leaves the scope because $b$ does not occur in it.

\begin{remark}
\label{claus:rem:e3pkgs}
In the Lean development the equivalences are
\Lean{e3prf\_iff\_vprf} on every clause context and
\Lean{e3prf\_iff\_nprf\_of\_qfree}. The Lean
translation from $\mathsf{V}$ carries a renaming and an extension of the
context through its induction, where the proof above uses
\cref{claus:lem:theta}. The translation of \rulename{Elim} copies the
derivation of its major premise into the scope of the witness.
\end{remark}

\begin{remark}[Checking without fuel]
\label{claus:rem:fuel}
\pkg{evaluation\_discharge} observes that the certificate predicate of
$\mathsf{ED}$ is primitive recursive with no fuel field. On tree sizes
raw substitution satisfies
$|a\subst{t}{u}|\le|a|\cdot\max(1,|t|)$, a step of $\varrho$ performs at
most one substitution in a body, each recursive call consumes a proper
subcertificate, and the work at a node is bounded by a fixed polynomial
in the sizes of its data, its context and the conclusions of its
subderivations. The zero-normal companion adds structural deletions and
at most one constructed iterator.
\end{remark}

\subsection{The zero-normal companion}

\begin{definition}[Seed evaluation and the calculus $\mathsf{ED}^{0}$]
\label{claus:def:e30}
On zero-normal terms let
\[
 \sigma_E(F)\defeq
 \begin{cases}
 r=h\bigl\langle\overline{I}(c,t;\,u.\,h)/u\bigr\rangle,
   &E\equiv(s=t),\ F\equiv\bigl(\It{\Suc c}{s}{u}{h}=r\bigr),\\
 F,&\text{otherwise.}
 \end{cases}
\]
$\mathsf{ED}^{0}$ has zero-normal terms, normal substitution, uniform
clauses, \rulename{Use} and \rulename{Close} of $\mathsf{V}^{0}$, and
\rulename{Discharge} with $\sigma_E(F)$ in the place of $\varrho_E(F)$.
The calculus $\mathsf{V}^{0}_{\mathrm{s}}$ is $\mathsf{V}^{0}$ with
\rulename{Step} replaced by the rule \rulename{SeedEval} of
\pkg{seed\_evaluation} (\cref{comp:def:seedeval}): from $s=t$ and
$\It{\Suc c}{s}{u}{h}=r$ infer
$r=h\langle\overline{I}(c,t;\,u.\,h)/u\rangle$.
\end{definition}

The case of $\sigma_E(F)$ that is not the identity is exactly
\rulename{SeedEval}, with $E$ and $F$ as its two premises.

\begin{lemma}[\rulename{SeedEval} and \rulename{Step}]
\label{claus:lem:seed}
On zero-normal terms, in every context that contains $\mathcal{R}$,
\rulename{SeedEval} is derivable from \rulename{Use} of $\mathcal{R}$ and
\rulename{Step}, and \rulename{Step} is derivable from \rulename{Use} of
$\mathcal{R}$ and \rulename{SeedEval}. No logical rule is used.
\end{lemma}
\begin{proof}
These are parts (b) and (c) of \cref{comp:lem:seedstep}. Their derivations
consist of reflexivity, \rulename{Step} or \rulename{SeedEval}, and the
given premises; they read no clause, so they are derivations over every
clause context once each reflexivity leaf is replaced by \rulename{Use} of
$\mathcal{R}$ at the term in question.
\end{proof}

\begin{theorem}[The zero-normal fusion]
\label{claus:thm:fusionzero}
For every zero-normal context $\Gamma$ of uniform clauses and every
zero-normal formula $A$,
$\Gamma\proves_{\mathsf{ED}^{0}}A$ if and only if
$\Gamma\proves_{\mathsf{V}^{0}_{\mathrm{s}}}A$. If $\Gamma$ contains
$\mathcal{R}$, both are equivalent to $\Gamma\proves_{\mathsf{V}^{0}}A$.
\end{theorem}
\begin{proof}
The proof of \cref{claus:thm:fusion} applies with \rulename{SeedEval} in
the place of the two rows of \rulename{Eval}: a major premise with root
$D$ recovers \rulename{Open}, and the major premise $\Theta$ recovers
\rulename{SeedEval}. No substitution in a derivation is used. The second
statement is \cref{claus:lem:seed}.
\end{proof}

The companion and its equivalence are not formalized in Lean; the
formalized calculus with \rulename{SeedEval} has native contexts
(\Lean{NVPrf} in \Lean{NativeZeroSeed}, \cref{comp:thm:seedeval}). With clause contexts and a rule
\rulename{Refl}, \rulename{SeedEval} gives the calculus
$\mathsf{C}_{\mathrm{Seed}}$ of \cref{comp:thm:seedeval}, which is
equivalent to $\mathsf{C}^{0}$.

\subsection{Partial evaluation and the diagonal introduction}
\label{claus:sec:pd}

This subsection draws on \pkg{476\_byte\_presentation}. Two rules of
$\mathsf{ED}$ have instances that a substitution does not carry to
instances of the same rule: the introduction selects its mode by the
witness (\cref{claus:prop:subst}), and the identity case of $\varrho$
applies to every pair that matches no arithmetic case
(\cref{claus:rem:total}). The calculus $\mathsf{PD}$ below takes the
diagonal introduction of $\mathsf{K}$ and restricts the identity case to an
$E$ with root $D$. It keeps three rule families, and substitution preserves
its derivation trees.

\begin{definition}[Partial evaluation]
\label{claus:def:rhoplus}
For formulas $E$ and $F$ the formula $\varrho^{+}_E(F)$ is
\begin{align*}
 \varrho^{+}_{\,s=t}(b=r)&\defeq\bigl(r=\Phi_2(s,t,b)\bigr)
   \qquad\text{if }\Phi_2(s,t,b)\text{ is defined},\\
 \varrho^{+}_{\,\D{p}{q}C}(F)&\defeq F,
\end{align*}
and $\varrho^{+}_E(F)$ is undefined in every other case.
\end{definition}

Written out, the first line consists of the zero row and the successor row
\[
 \varrho^{+}_{\,s=t}\bigl(\It{0}{s}{u}{h}=r\bigr)\equiv(r=t),\qquad
 \varrho^{+}_{\,\It{c}{a}{u}{h}=t}\bigl(\It{\Suc c}{a}{u}{h}=r\bigr)
 \equiv\bigl(r=h\subst{t}{u}\bigr),
\]
which are the first two cases of $\varrho$ (\cref{claus:def:rho}); the
second line, the \emph{row} $D$, is the third case of $\varrho$ restricted
to an $E$ with root $D$. The three rows are disjoint: in the first two $E$
is an equation and the counter of the left side of $F$ is $0$,
respectively a successor, and in row $D$ the formula $E$ has root $D$.
Hence $\varrho^{+}$ is a partial function, and
$\varrho^{+}_E(F)\equiv\varrho_E(F)$ wherever $\varrho^{+}_E(F)$ is
defined. A pair of equations that matches no row of $\Phi_2$ is not in
its domain.

\begin{definition}[The calculi $\mathsf{K}^{u}$ and $\mathsf{PD}$]
\label{claus:def:pd}
$\mathsf{K}^{u}$ is the calculus $\mathsf{V}$ of \cref{claus:def:v} with
the term $\eta$ of \rulename{Intro} replaced by the term $\eta_B(z)$ of the
diagonal introduction (\cref{claus:def:k}); its rule families are
\rulename{Use}, \rulename{Eval}, \rulename{Elim} and \rulename{Intro}.
$\mathsf{PD}$ has raw terms, contexts of uniform clauses, the rules
\rulename{Use} and \rulename{Intro} of $\mathsf{K}^{u}$ and
\[
 \infrule{\Gamma\proves\D{x}{y}B\quad
          \Gamma,\mathcal{B}_b\proves E\quad
          \Gamma,\mathcal{B}_b\proves F}
         {\Gamma\proves\varrho^{+}_E(F)}\ \rulename{Discharge}
\]
where $\mathcal{B}_b\defeq[B(b,0);\,y.\,B(b,y)]$, the formula
$\varrho^{+}_E(F)$ is defined, and $b$ is fresh for $\Gamma$ and $B$ and
does not occur free in $\varrho^{+}_E(F)$.
\end{definition}

$\mathsf{K}^{u}$ stands to $\mathsf{K}$ as $\mathsf{V}$ stands to
$\mathsf{C}$. The formulas $E$ and $F$ of a \rulename{Discharge} may
contain $b$. The specification of $\mathsf{PD}$ is the clause
$\mathcal{R}$ in the initial context of a theorem, \rulename{Use} with zero
or one premise, \rulename{Intro} with the diagonal test made on the matrix
before the witness is substituted, and \rulename{Discharge} with the three
rows of $\varrho^{+}$; there is no reflexivity, computation, transport,
induction, cut or structural rule. Every \rulename{Discharge} of
$\mathsf{PD}$ is a \rulename{Discharge} of $\mathsf{ED}$, and the two
calculi differ in the mode of \rulename{Intro} and in the domain of the
identity case.

\begin{lemma}[Covariance of $\varrho^{+}$]
\label{claus:lem:covar}
Let $\sigma$ be a capture-free substitution of terms for free variables.
If $\varrho^{+}_E(F)$ is defined, then $\varrho^{+}_{E\sigma}(F\sigma)$ is
defined by the same row and
\[
 \varrho^{+}_{E\sigma}(F\sigma)\equiv\varrho^{+}_E(F)\,\sigma .
\]
\end{lemma}
\begin{proof}
In the zero and the successor row $E\equiv(s=t)$ and $F\equiv(b=r)$, so
$E\sigma\equiv(s\sigma=t\sigma)$ and $F\sigma\equiv(b\sigma=r\sigma)$. By
\cref{comp:prop:phisubst}, $\Phi_2(s\sigma,t\sigma,b\sigma)$ is defined by
the same row and is $\Phi_2(s,t,b)\sigma$; hence
$\varrho^{+}_{E\sigma}(F\sigma)\equiv(r\sigma=\Phi_2(s,t,b)\sigma)
\equiv\varrho^{+}_E(F)\sigma$. In row $D$ the formula $E\sigma$ has root
$D$ (\cref{guard:prop:subst-shapes}(i)), and both sides are $F\sigma$.
\end{proof}

Definedness is not reflected: the substitution $x\mapsto 0$ takes the pair
$(x=0,\ \It{0}{0}{u}{0}=0)$, outside the domain, to a pair in the zero row.
The theorems below need preservation only.

\begin{remark}[The total identity case is not covariant]
\label{claus:rem:total}
Let $E\equiv(x=0)$, $F\equiv(\It{0}{0}{u}{0}=0)$ and $\sigma(x)\defeq 0$.
No arithmetic case of $\varrho$ matches $(E,F)$: the zero case needs the
seed $0$ of the iterator to be the left side $x$ of $E$, and the successor
case needs a successor counter. So $\varrho_E(F)\equiv F$. The pair
$(E\sigma,F\sigma)\equiv(0=0,\ \It{0}{0}{u}{0}=0)$ matches the zero case,
so
\[
 \varrho_{E\sigma}(F\sigma)\equiv(0=0)\not\equiv
 \bigl(\It{0}{0}{u}{0}=0\bigr)\equiv\varrho_E(F)\,\sigma .
\]
A \rulename{Discharge} of $\mathsf{ED}$ at this pair therefore does not
become, under $\sigma$, an instance of the same rule with the substituted
conclusion. This obstructs substitution that preserves the derivation tree
of $\mathsf{ED}$, not admissibility: for contexts $\mathcal{R},\Gamma$,
substitution is admissible in $\mathsf{ED}$ through
\cref{claus:thm:fusion}, \cref{claus:thm:uniform} and
\cref{claus:prop:subst}, since $(\Gamma\sigma)^{\circ}\equiv
\Gamma^{\circ}\sigma$. The pair is outside the domain of $\varrho^{+}$.
\end{remark}

For a substitution $\sigma$, the image of a uniform clause
$[\Pi;\,y.\,C]$ is $[\Pi\sigma;\,y.\,C\sigma]$, with $y$ renamed away from
the variables and values of $\sigma$; then
$(C\subst{t}{y})\sigma\equiv(C\sigma)\subst{t\sigma}{y}$,
$\langle A\rangle\sigma\equiv\langle A\sigma\rangle$ and
$\mathcal{R}\sigma\equiv\mathcal{R}$.

\begin{theorem}[Substitution and weakening preserve the derivation tree]
\label{claus:thm:pdsubst}
Let $\Gamma\proves_{\mathsf{PD}}A$, let $\sigma$ be a capture-free
substitution of terms for free variables, and let $\Gamma'$ be a context
with a map $\iota$ from the occurrences of $\Gamma$ to the occurrences of
$\Gamma'$ such that $\iota(c)$ is the clause $c\sigma$. Then
$\Gamma'\proves_{\mathsf{PD}}A\sigma$ by a derivation with the same tree
of rules, in which every \rulename{Use} selects the image of its
occurrence, every \rulename{Intro} has the same mode and every
\rulename{Discharge} the same row of $\varrho^{+}$, after a renaming of
eigenvariables. In particular substitution ($\Gamma'\defeq\Gamma\sigma$),
weakening and exchange ($\sigma$ the identity) and injective renaming are
admissible in $\mathsf{PD}$.
\end{theorem}
\begin{proof}
Induction on the derivation, simultaneously for all $\sigma$, $\Gamma'$
and $\iota$.

\rulename{Use}: the derivation selects $c\equiv[\Pi;\,y.\,C]$ at the term
$t$ and has a derivation of $\Gamma\proves A'$ for each $A'$ in $\Pi$. The
induction hypothesis gives $\Gamma'\proves A'\sigma$ for each member of
$\Pi\sigma$, and \rulename{Use} of $\iota(c)\equiv[\Pi\sigma;\,y.\,C\sigma]$
at $t\sigma$ gives $(C\sigma)\subst{t\sigma}{y}\equiv(C\subst{t}{y})\sigma$.

\rulename{Intro}: the premise is
$\Gamma,\langle B(t,\eta_B(z))\rangle\proves B(t,\Suc z)$ with $z$ fresh
for $\Gamma$, $t$ and $B$. Choose $z'$ fresh for $\Gamma'$, $t\sigma$ and
$B\sigma$, and let $\sigma'$ be $\sigma$ with $z\mapsto z'$ in addition.
Since $z$ does not occur in $\Gamma$, $t$ and $B$, the substitution
$\sigma'$ acts on them as $\sigma$, so $\iota$ still satisfies the
hypothesis for $\sigma'$. The diagonal test is preserved and reflected by
the substitution, lifted under $x$ and $y$
(\cref{guard:prop:subst-shapes}(iii)), so
$\eta_B(z)\sigma'\equiv\eta_{B\sigma}(z')$ and the added occurrence has the
image $\langle(B\sigma)(t\sigma,\eta_{B\sigma}(z'))\rangle$. Extend
$\iota$ by sending the added occurrence to the occurrence added on the
right. The induction hypothesis for $\sigma'$ gives
$\Gamma',\langle(B\sigma)(t\sigma,\eta_{B\sigma}(z'))\rangle\proves
(B\sigma)(t\sigma,\Suc z')$, and \rulename{Intro} in the same mode with
the eigenvariable $z'$ gives $\Gamma'\proves\D{x}{y}(B\sigma)\equiv
(\D{x}{y}B)\sigma$.

\rulename{Discharge}: the premises are $\Gamma\proves\D{x}{y}B$ and
$\Gamma,\mathcal{B}_b\proves E$, $\Gamma,\mathcal{B}_b\proves F$, with
$\varrho^{+}_E(F)\equiv R$ and $b$ fresh for $\Gamma$, $B$ and $R$. The
induction hypothesis for $\sigma$ and $\iota$ gives
$\Gamma'\proves\D{x}{y}(B\sigma)$. Choose $b'$ fresh for $\Gamma'$,
$B\sigma$ and $R\sigma$, and let $\sigma'$ be $\sigma$ with $b\mapsto b'$
in addition. Then $\mathcal{B}_b\sigma'$ is the clause
$[(B\sigma)(b',0);\,y.\,(B\sigma)(b',y)]$ that opens $\D{x}{y}(B\sigma)$
at $b'$; extend $\iota$ accordingly. The induction hypothesis for
$\sigma'$ gives the two branches with the conclusions $E\sigma'$ and
$F\sigma'$. By \cref{claus:lem:covar},
$\varrho^{+}_{E\sigma'}(F\sigma')$ is defined by the same row and is
$R\sigma'\equiv R\sigma$, because $b$ does not occur in $R$. The
eigenvariable condition holds for $b'$, and \rulename{Discharge} gives
$\Gamma'\proves R\sigma$.

For the special cases take $\iota$ positional for $\Gamma\sigma$, and for
weakening and exchange the map that sends each occurrence of $\Gamma$ to
its copy in $\Gamma'$.
% ed.: the package renames the eigenvariables first and then applies the
% induction hypothesis to the renamed branch; extending sigma by z -> z'
% and b -> b' applies the hypothesis to the branch itself, so no separate
% renaming lemma is presupposed.
\end{proof}

No cut, normalization or arithmetic bridge is used. Substitution in
$\mathsf{K}$ preserves the derivation tree as well
(\cref{claus:thm:ksubst}); $\mathsf{K}$ has four rule families and an
explicit \rulename{Refl}, and $\mathsf{PD}$ has three and the clause
$\mathcal{R}$.

\begin{theorem}[Uniform clauses with the diagonal introduction]
\label{claus:thm:ku}
For every context $\Gamma$ of uniform clauses and every formula $A$,
\[
 \mathcal{R},\Gamma\proves_{\mathsf{K}^{u}}A\quad\Longleftrightarrow\quad
 \Gamma^{\circ}\proves_{\mathsf{K}}A .
\]
Hence, for $Q$-free $\Delta$ and $F$,
$\mathcal{R},\langle\Delta\rangle\proves_{\mathsf{K}^{u}}F$ if and only if
$\Delta\proves_{\Nat}F$.
\end{theorem}
\begin{proof}
The proof of \cref{claus:thm:uniform} applies. Its invariant is that each
occurrence of the uniform context, other than $\mathcal{R}$, corresponds by
position to an occurrence of the packed context with the same consequent
and either the same antecedent or no antecedent in the place of $0=0$,
and that \rulename{Use} of $\mathcal{R}$ corresponds to \rulename{Refl}.
It copies every \rulename{Intro} with the same matrix,
witness and eigenvariable, and the added occurrences
$\langle B(t,\eta)\rangle$ and $\operatorname{ord}(B(t,\eta))\equiv
\langle B(t,\eta)\rangle^{\circ}$ correspond. Here $\eta\equiv\eta_B(z)$
on both sides, since the diagonal test reads the matrix only. An
\rulename{Elim} adds the same opened occurrence on both sides, which has an
antecedent. The second statement follows from \cref{claus:thm:k}, that is
statement (iii) of \cref{claus:thm:five} for $\mathsf{K}$, since
$\langle\Delta\rangle^{\circ}\equiv\operatorname{ord}(\Delta)$.
\end{proof}

\begin{theorem}[Partial evaluation at the discharge]
\label{claus:thm:pd}
For every context $\Gamma$ of uniform clauses and every formula $A$,
\[
 \Gamma\proves_{\mathsf{PD}}A\quad\Longleftrightarrow\quad
 \Gamma\proves_{\mathsf{K}^{u}}A .
\]
The translations leave the end sequent unchanged, and $\mathcal{R}$ need
not be in $\Gamma$. Hence, for $Q$-free $\Delta$ and $F$,
$\mathcal{R},\langle\Delta\rangle\proves_{\mathsf{PD}}F$ if and only if
$\Delta\proves_{\Nat}F$.
\end{theorem}
\begin{proof}
Both directions are inductions on derivations. \rulename{Use} and
\rulename{Intro} are the same rules in both calculi and are copied.

From left to right translate the three premises of a \rulename{Discharge}
with major premise $\D{x}{y}B$ and conclusion $R\equiv\varrho^{+}_E(F)$.
In the zero or the successor row, $E\equiv(s=t)$, $F\equiv(b=r)$ and
$R\equiv(r=\Phi_2(s,t,b))$; \rulename{Eval} applied to the translated
derivations of $E$ and $F$ proves $R$ in $\Gamma,\mathcal{B}_b$. In row $D$,
$R\equiv F$, and the translated derivation of $F$ proves it. Then apply
\rulename{Elim} to the translated major premise; its eigenvariable
condition is the condition of \rulename{Discharge} on $b$.

From right to left let $\Theta\equiv\D{x}{y}(0=0)$. In $\mathsf{PD}$ it is
derivable in every context, also without $\mathcal{R}$, by the two-node
derivation of \cref{claus:lem:theta}: the matrix $0=0$ is not diagonal, so
\rulename{Intro} at the witness $0$ is in case mode, its assumption and its
target are both $0=0$, and \rulename{Use} of $\langle 0=0\rangle$ closes the
branch. Translate an \rulename{Elim} with major premise $\D{x}{y}B$ and
branch $\Gamma,\mathcal{B}_b\proves A$ by \rulename{Discharge} with the
translated major premise, with $E\defeq\Theta$, derived in
$\Gamma,\mathcal{B}_b$, and with $F\defeq A$; row $D$ gives $A$, and the
eigenvariable condition is that of \rulename{Elim}. Translate an
\rulename{Eval} with premises $\Gamma\proves s=t$ and $\Gamma\proves b=r$
by \rulename{Discharge} with the major premise $\Theta$, opened at a $b_0$
fresh for $\Gamma$ and the four terms, which adds the occurrence
$[0=0;\,y.\,0=0]$. Weaken the translated derivations of the two premises
into $\Gamma,[0=0;\,y.\,0=0]$ by \cref{claus:thm:pdsubst}. The pair
$(s=t,\ b=r)$ is in the zero or the successor row of $\varrho^{+}$ with the
value $r=\Phi_2(s,t,b)$, the conclusion of the \rulename{Eval}, in which
$b_0$ does not occur.

The second statement is \cref{claus:thm:ku}.
\end{proof}

\begin{remark}[The Lean modules]
\label{claus:rem:pdlean}
\Lean{NativePartialDischarge} formalizes $\varrho^{+}$ (\Lean{RhoPlus}, with
$\varrho^{+}\subseteq\varrho$ as \Lean{RhoPlus.toRho}), $\mathsf{K}^{u}$
(\Lean{KUPrf}) and $\mathsf{PD}$ (\Lean{PDPrf}), and proves
\cref{claus:thm:pd} for every clause context (\Lean{pdprf\_iff\_kuprf}),
\cref{claus:thm:ku} with the reflexivity clause (\Lean{kuprf\_iff\_kprf}),
and the bridge on $Q$-free sequents (\Lean{pdprf\_iff\_nprf\_of\_qfree}).
The proofs follow those of the evaluating discharge (\Lean{E3Prf}) and of
uniform clauses (\Lean{VPrf}): the diagonal start of the introduction
depends on the matrix only and commutes with substitution
(\Lean{NFm.dstart\_subst}), and the row $D$ takes the place of the identity
case.  The translation from $\mathsf{K}^{u}$ carries an injective renaming
and an extension of the context through the induction instead of a
separate substitution lemma, so \cref{claus:thm:pdsubst} is not needed
for it and is not formalized.
\end{remark}

Under the ordinary reading every rule of $\mathsf{PD}$ is sound in $\Teq$
directly: \rulename{Discharge} as the discharge of $\mathsf{ED}$ after
\cref{claus:thm:fusion}, since $\varrho^{+}$ is a restriction of $\varrho$,
and the diagonal \rulename{Intro} by coverage in case mode and by induction
for the literal equation under the assumed base in progress mode. By
\cref{claus:thm:pd}, $\mathsf{PD}$ and $\mathsf{K}^{u}$ are among the
calculi of \cref{claus:cor:bridge}: through \Lean{kprf\_iff\_nprf\_of\_qfree}
and \cref{base:thm:bridge} they are bi-interpretable with $\iPRA$ with
whole domains, actual equality and identity numerical composites, and
$\mathcal{R}\proves_{\mathsf{PD}}0=\Suc 0$ if and only if $\iPRA$ derives
falsity.

\begin{remark}
\label{claus:rem:pdpkgs}
\pkg{476\_byte\_presentation} states \cref{claus:thm:pdsubst} for
substitution and, without proof details, for weakening and renaming, and
proves \cref{claus:thm:pd} through $\mathsf{K}^{u}$ for contexts that
contain $\mathcal{R}$; the proof above does not use $\mathcal{R}$. Its
derivation of $\Theta$ is that of \cref{claus:lem:theta}, its derivation of
the standalone evaluation from \rulename{Discharge} at $\Theta$ is that of
the proof of \cref{claus:thm:fusion}, and its translation of
\rulename{Discharge} into \rulename{Eval} and \rulename{Elim} is the
converse direction of that proof, restricted to the domain of
$\varrho^{+}$. The one difference is the translation of \rulename{Elim}:
the proof of \cref{claus:thm:fusion} takes as $E$ the major premise
weakened into the scope of the witness, and the proof above takes the
two-node derivation of $\Theta$, which does not copy the derivation of the
major premise. Both choices are available in either calculus, since the
weakened major premise also has root $D$. The package also recovers
reflexivity by \rulename{Use} of $\mathcal{R}$ and, from the derived
standalone evaluation, the two computation equations, the inference
\rulename{Root} by the frame $\It{0}{s}{u}{u}$, symmetry, transitivity and
one-template transport; these are the derivations of \cref{claus:lem:eq}.
The package compiles certificates in both directions between its checker
and the matched four-family checker of $\mathsf{K}^{u}$, preserving the
goal; the compilers are untrusted and their outputs are rechecked.
\end{remark}

\section{The ordinary realization}
\label{claus:sec:real}

This section draws on all twelve packages of the chapter:
\pkg{raw\_clausal}, \pkg{rootfree\_clausal},
\pkg{reduced\_clause\_calculus}, \pkg{clausal\_two\_case},
\pkg{clausal\_single\_evaluation}, \pkg{uniform\_clauses},
\pkg{clause\_guards}, \pkg{clause\_contexts}, \pkg{prefix\_clauses},
\pkg{one\_body}, \pkg{evaluation\_discharge} and
\pkg{476\_byte\_presentation}. Each restates the
ordinary realization, the logical macros of the native interface (gate,
selector and parity term, with $q_y$ in the place of $\Q{y}$) and the
bridge. These belong to the common core, and the statements reduce to it
as follows.

\begin{corollary}[Bi-interpretability]
\label{claus:cor:bridge}
Let $\mathsf{X}$ be one of $\mathsf{C}_3$, $\mathsf{C}$, $\mathsf{C}^{0}$,
$\mathsf{C}^{0}_{\mathrm{b}}$, $\mathsf{V}$, $\mathsf{V}^{0}$,
$\mathsf{K}$, the mirror calculus, $\mathsf{PF}$, $\mathsf{OB}$,
$\mathsf{HB}$, $\mathsf{ED}$, $\mathsf{ED}^{0}$, $\mathsf{K}^{u}$,
$\mathsf{PD}$, with its fixed
reflexivity clause or entry where it has one. Then $\mathsf{X}$ derives
$0=\Suc 0$ from no assumptions if and only if $\Nat$ does, if and only if
$\iPRA$ derives falsity. The ordinary realization of $\mathsf{X}$ is
bi-interpretable with $\iPRA$ with whole domains, actual equality and
identity numerical composites.
\end{corollary}
\begin{proof}
Each calculus has been shown equivalent to $\Nat$ on $Q$-free sequents by
translations that are the identity on terms, or the map $\nf_0$ on terms
for the zero-normal calculi, and the formula composites are the identity
and $A\mapsto j\tau A$, respectively $A\mapsto\nf_0(j\tau A)$, which are
provably equivalent to $A$ in $\Nat$
(\cref{claus:lem:unq,comp:lem:nusound}). The claim follows from
\cref{base:thm:bridge}.
\end{proof}

Under the ordinary reading every rule is sound in $\Teq$ directly:
\rulename{Refl}, \rulename{Eval}, \rulename{Step} and \rulename{SeedEval}
follow from the computation equations and equality, \rulename{Use} and
the eliminations are ordinary eliminations, an introduction in case mode
uses coverage, and an introduction in progress mode uses induction for a
literal equation under the assumed base. The separation
$0=\Suc 0\to\bot$ is an axiom of the ordinary realization and not a rule
of a clause calculus, in which $0=\Suc 0$ is the designated contradiction;
positive equations alone admit a one-element structure.

\section{Successor rigidity}
\label{claus:sec:rigidity}

This section draws on \pkg{raw\_clausal}. It concerns an arithmetic
principle and not a clause calculus. The background theory is
intuitionistic first-order logic with equality over the scalar terms,
with the two computation equations and coverage
$\forall n\,(n=0\lor\exists m\,n=\Suc m)$. Neither the injectivity of
$\Suc$ nor the separation of $0$ from $\Suc 0$ is used. The rule is
\[
 \infrule{\Gamma\proves f(0)=0\quad\Gamma\proves f(\Suc z)=\Suc(f(z))}
         {\Gamma\proves f(t)=t}\ \mathrm{SR}_0
\]
where $f(z)$ is a term with a distinguished variable, not a function
variable, and $z$ is fresh for $\Gamma$, $t$ and the parameters of $f$.
The second premise has no local assumption. The rule is the uniqueness of
iteration \rulename{UR} (\cref{kern:sec:ur}) restricted to the initializer
$0$ and the update $\Suc$. Induction for literal equations derives
$\mathrm{SR}_0$: induct on $f(n)=n$, where the step is
$f(\Suc z)=\Suc(f(z))=\Suc z$ under $f(z)=z$. Successor uniqueness
\rulename{SU} (\cref{kern:def:SU}) also derives it: \rulename{SU} gives
$f(t)=\It{t}{f(0)}{u}{\Suc u}$, the first premise turns the seed into $0$,
and $\It{t}{0}{u}{\Suc u}=t$ is \cref{kern:lem:cu}(iii), which is
derivable in $T_{\mathrm{SU}}$.

\begin{theorem}[Zero-equation induction from successor rigidity]
\label{claus:thm:rigidity}
Work in the background theory with $\mathrm{SR}_0$. Let $E(x)$ be a
formula with $\Gamma\proves E(0)$ and $\Gamma,E(z)\proves E(\Suc z)$, $z$
fresh, and let $e(x)$ be a term such that, for a fresh $x$,
\[
 \Gamma\proves E(e(x)),\qquad \Gamma,E(x)\proves e(x)=x .
\]
Then $\Gamma\proves E(t)$ for every term $t$. In particular, for every
term $q(x)$ the rule
\[
 \infrule{\Gamma\proves q(0)=0\quad\Gamma,q(z)=0\proves q(\Suc z)=0}
         {\Gamma\proves q(t)=0}\ \mathrm{ZI}
\]
is derivable.
\end{theorem}
\begin{proof}
Let $r(n)\defeq\It{n}{0}{u}{e(\Suc u)}$. First, $\Gamma\proves E(r(n))$,
by coverage on $n$. If $n=0$ then $r(n)=0$ and the base applies. If
$n=\Suc m$ then $r(n)=e(\Suc(r(m)))$, and the first condition on $e$ at
$\Suc(r(m))$ applies. This step assumes nothing about $r(m)$ and uses no
induction. Next instantiate the given step at $r(z)$; with $E(r(z))$ it
gives $E(\Suc(r(z)))$, and the second condition on $e$ gives
$e(\Suc(r(z)))=\Suc(r(z))$. The computation equations therefore prove
$r(0)=0$ and $r(\Suc z)=\Suc(r(z))$ in the context $\Gamma$.
$\mathrm{SR}_0$ gives $r(t)=t$, and $E(r(t))$ transports to $E(t)$. The
eigenvariable of $\mathrm{SR}_0$ and the witness of the coverage
elimination are fresh, so $\Gamma$ may contain arbitrary formulas.

For $\mathrm{ZI}$ take $E(x)\equiv(q(x)=0)$ and
$e(x)\defeq\sel(q(x),x,0)$. If $q(x)=0$ the selector computes to $x$,
which is the second condition. For the first apply coverage to $q(x)$. If
$q(x)=0$ then $e(x)=x$ and $q(e(x))=0$. If $q(x)=\Suc v$ then $e(x)=0$
and the base $q(0)=0$ gives $q(e(x))=0$. These cases need no test of
equality of two terms and no separation.
\end{proof}

The construction of $e$ and $r$ is effective on syntax, so the theorem
is a transformation of derivations. The theorem does not give full
scalar arithmetic from $\mathrm{SR}_0$; see \cref{claus:q:rigidity}.

\begin{remark}[Comparison with the rigidity of retractions]
\label{claus:rem:rigidity}
\Cref{kern:prop:rinv,kern:lem:rigid,kern:lem:split,kern:lem:zi,kern:thm:su}
derive induction for zero equations from \rulename{SU} by a trajectory of
the same kind, $r(n)\equiv\It{n}{0}{u}{\Suc e(u)}$, for a term $e$ with
$e(0)=0$ and $e(\Suc e(x))=\Suc e(x)$; since \pkg{reductions} that
derivation needs neither idempotence of $e$ nor coverage
(\cref{kern:prop:rinv}), and coverage and separation enter only in
\cref{kern:lem:zi}. Since $\mathrm{SR}_0$ is derivable from
\rulename{SU}, \cref{claus:thm:rigidity} implies \cref{kern:thm:su}.
What it adds is the following. Its rule $\mathrm{SR}_0$ fixes the seed $0$
and concludes the identity, so it does not presuppose
\cref{kern:lem:cu}(iii). Its two conditions on $e$ concern the formula $E$
and not $e$ alone, so the general statement applies to every formula $E$
with such an $e$. Its derivation of $\mathrm{ZI}$ uses coverage on $q(x)$
in place of a decided split (\cref{kern:lem:split}), so it needs neither
the decidability of $q(n)=0$ nor separation, which \cref{kern:lem:zi}
uses.
\end{remark}

\section{Checker sizes}
\label{claus:sec:checkers}

This section collects the measurements of the twelve packages named in
\cref{claus:sec:real}. All of them measure complete standalone Python
checkers by the protocol of \cref{base:sec:measure}, under CPython
3.12.14: the module without its
docstring, printed by \texttt{ast.unparse}, with one final line feed.
Each count includes parsing, validation of syntax and scope, substitution,
the rules and the acceptance interface. The checkers read one JSON value,
a closed goal and a proof, and accept closed theorems; open sequents occur
in the branches. The universal statements of this chapter are the written
proofs; the packages' finite test suites concern the implementations.

\emph{Representation.} With three exceptions the checkers share one
codec:
$0$ is the integer $0$, a free parameter is a negative level, a bound
variable is a positive index, $\Suc t$ is \texttt{[0,t]}, the iterator is
\texttt{[n,a,h]}, an equation is \texttt{[s,t]} and $\D{x}{y}B$ is
\texttt{[B]}, the expected sort separating a successor from an equation.
Every inserted term has no loose positive index, so opening a binder needs
no shifting; an elimination or introduction assigns one new negative level
and validates its conclusion in the outer scope, which enforces the
eigenvariable condition. \pkg{clausal\_two\_case} uses tagged lists and
infers the conclusion of an elimination; it proves, by induction on a
successful computation, that a returned formula is well scoped and
derivable in $\mathsf{C}$ from the retained context, and that every
derivation of a closed theorem has a certificate. \pkg{clause\_contexts}
writes $\Suc t$ as \texttt{[t]} and $\D{x}{y}B$ as \texttt{["d",B]}.
\pkg{prefix\_clauses} stores a formula as \texttt{[k,s,t]} and $\Suc t$ as
\texttt{[t]}; the expected sort separates a record from an iterator.

\emph{Protocols.} \pkg{clausal\_single\_evaluation} and
\pkg{clause\_contexts} check goal-directed certificates in which a
validated reflexive goal closes without a proof node. The selected
checker of \pkg{clause\_contexts} reads a JSON array of two-field packets
followed by the goal, from right to left; the shape of the first field
selects \rulename{Use}, \rulename{Close}, \rulename{Open} or
\rulename{Step}, an unconsumed prefix is rejected, and an unconditional
local clause is stored with the antecedent $0=0$, which closes without
input. The package proves that a successful check yields a derivation of
$\mathsf{V}^{0}$ and that every derivation of a closed theorem has an
accepted certificate, by induction on the call tree and by the invariant
that a call consumes exactly its own segment. In \pkg{prefix\_clauses}
and \pkg{evaluation\_discharge} the premise field of a \rulename{Use}
whose premise is omitted or absent is an inert field of the format.
\pkg{evaluation\_discharge} infers every intermediate conclusion from the
proof tree and supplies no template to \rulename{Discharge}.

\emph{What the measurements show.} The deletion of the root row saves 47
bytes, 5 tokens and 7 AST nodes in an unchanged format. Uniform clauses
save 80 bytes against the same validator. Zero-normal syntax with one
evaluation scheme is smaller than its raw control in
\pkg{clausal\_single\_evaluation} and \pkg{rootfree\_clausal} and larger
in the companion of \pkg{reduced\_clause\_calculus}. The diagonal checker
is 100 bytes larger than the witness-mode checker of
\pkg{rootfree\_clausal}. In \pkg{clause\_contexts} the factorization of
clauses alone saves 53 bytes, 2 tokens and 5 AST nodes against the
strongest control, and the rest of its reduction comes from the packet
protocol and shared validation; its ranked representation is larger than
its recursive one. In \pkg{prefix\_clauses} the prefix representation
alone saves 3 tokens and 63 bytes and adds 10 AST nodes against its
nested control. The one-body checker and the homogeneous checker are
larger than their controls, by the cost of the gate and its lifting. The
fusion at the discharge saves 103 bytes, 22 tokens and 22 AST nodes on raw
terms and 21 bytes, 12 tokens and 10 AST nodes on zero-normal terms.
\pkg{one\_body} also reports a refinement of an existing checker for the
native calculus on the equation and $D$ language that changes no rule.
The checker of $\mathsf{PD}$ in \pkg{476\_byte\_presentation} measures
678 tokens, 727 AST nodes and 2087 bytes, against 688/735/2128 for its
matched split control (the four families of $\mathsf{K}^{u}$, same strict
certificate contract) and 643/695/1892 for the inherited checker of
$\mathsf{ED}$, so it saves 10 tokens, 8 AST nodes and 41 bytes against its
control and sets no size record.

\begin{longtable}{@{}>{\raggedright\arraybackslash}p{0.15\linewidth}>{\raggedright\arraybackslash}p{0.23\linewidth}rrr>{\raggedright\arraybackslash}p{0.30\linewidth}@{}}
\caption{Checker sizes as stated by the packages: tokens, AST nodes,
normalized bytes. A control is given as tokens/AST nodes/bytes.}
\label{claus:tab:sizes}\\
\toprule
Package & Program & Tok. & AST & Bytes & Matched control\\
\midrule
\endfirsthead
\toprule
Package & Program & Tok. & AST & Bytes & Matched control\\
\midrule
\endhead
\bottomrule
\endfoot
\pkg{raw\_clausal} & five rules, three rows ($\mathsf{C}_3$) & 853 & 993 &
3046 & canonical five-rule control 903/1035/3175; raw six-rule control
930/1072/3252; upstream canonical checker 946/1072/3305\\
\pkg{rootfree\_\allowbreak clausal} &raw, two rows ($\mathsf{C}$) & 837 & 977 & 2956 &
matched raw control 842/984/3003; incoming raw clausal 853/993/3046\\
\pkg{rootfree\_\allowbreak clausal} &zero-normal, fixed zero base
($\mathsf{C}^{0}_{\mathrm{b}}$) & 852 & 987 & 2948 & zero-normal control
with root row 864/1006/2993\\
\pkg{rootfree\_\allowbreak clausal} &zero-normal, witness test ($\mathsf{C}^{0}$) &
853 & 992 & 2948 & the same control; tagged native normal comparator
979/1047/3011\\
\pkg{reduced\_clause\_\allowbreak calculus} & raw, two rows ($\mathsf{C}$)
& 848 & 986 & 2999 & exact raw clause control 853/993/3046\\
\pkg{reduced\_clause\_\allowbreak calculus} & companion, one scheme
($\mathsf{C}^{0}$) & 871 & 1011 & 3050 & zero-normal with root row
879/1023/3107\\
\pkg{clausal\_\allowbreak two\_case} &two-case checker ($\mathsf{C}$) & 794 & 868 &
2787 & matched three-case control 808/888/2861; pinned three-case checker
853/993/3046; the same with the root row deleted 848/986/2999\\
\pkg{clausal\_single\_\allowbreak evaluation} & single evaluation
($\mathsf{C}^{0}$) & 741 & 826 & 2283 & matched raw control 744/842/2416;
zero-normal control with root row 765/865/2426\\
\pkg{clausal\_single\_\allowbreak evaluation} & forward implementation &
867 & 1008 & 3022 & forward control 853/993/3046\\
\pkg{uniform\_\allowbreak clauses} &uniform clauses ($\mathsf{V}$) & 807 & 931 & 2912
& raw clause checker 853/993/3046; with shared validator 825/958/3039;
two rows 848/986/2999; two rows, shared validator 820/951/2992\\
\pkg{clause\_\allowbreak guards} &diagonal ($\mathsf{K}$) & 872 & 1013 & 3056 &
native diagonal control 975/1138/3395; witness-mode clause control
837/977/2956\\
\pkg{clause\_\allowbreak guards} &mirror guard & 874 & 1012 & 3059 & the same\\
\pkg{clause\_\allowbreak contexts} &four schemes, packets ($\mathsf{V}^{0}$) & 666 &
770 & 2282 & pinned five-pattern tagged source 829/902/2950; five-pattern
arity and scope control 753/828/2730\\
\pkg{clause\_\allowbreak contexts} &packets, empty antecedent & 665 & 770 & 2308 &
the same\\
\pkg{clause\_\allowbreak contexts} &explicit tree checker & 751 & 823 & 2677 & the
same\\
\pkg{clause\_\allowbreak contexts} &ranked tree checker & 781 & 853 & 2692 & the
same\\
\pkg{clause\_\allowbreak contexts} &support-rank experiment, not selected & 854 &
957 & 3069 & none\\
\pkg{prefix\_\allowbreak clauses} &prefix calculus ($\mathsf{PF}$) & 788 & 900 & 2700
& matched nested control 791/890/2763; \pkg{uniform\_clauses}
807/931/2912\\
\pkg{prefix\_\allowbreak clauses} &homogeneous variant ($\mathsf{HB}$) & 840 & 965 &
2843 & the same\\
\pkg{one\_body} & one body ($\mathsf{OB}$) & 901 & 1041 & 3130 & matched
two-field control 848/986/2999\\
\pkg{one\_body} & native checker, preferred refinement & 357 & 412 & 778 &
existing native checker 358/413/788\\
\pkg{one\_body} & native checker, byte-priority variant & 357 & 415 & 777
& the same\\
\pkg{evaluation\_\allowbreak discharge} & raw fusion ($\mathsf{ED}$) &
643 & 695 & 1892 & raw separate rules 665/717/1995\\
\pkg{evaluation\_\allowbreak discharge} & zero-normal fusion
($\mathsf{ED}^{0}$) & 633 & 697 & 1858 & zero-normal separate rules
645/707/1879; incoming controls 663/751/1983, 683/771/1990 and
636/708/2010\\
\end{longtable}

\section{Further questions}
\label{claus:sec:questions}

\begin{question}[Successor rigidity]
\label{claus:q:rigidity}
Do the computation equations, coverage, separation and $\mathrm{SR}_0$
construct a term $\delta(a,b)$ with $\delta(a,b)=0\leftrightarrow a=b$ on
the whole domain without using induction for arbitrary literal equations?
Such a term would turn $\mathrm{ZI}$ into induction for every literal
equation, and the bridge (\cref{base:thm:bridge}) would apply. The known
proofs of the predecessor (\cref{base:thm:pred}) and of the reflection of
equality (\cref{base:thm:reflect}) use inductions on general equations or
uniqueness with updates other than $\Suc$
(\crefrange{kern:sec:urpred}{kern:sec:urreflect}), so importing them is
circular.
\pkg{raw\_clausal} claims no converse from $\mathrm{ZI}$ to
$\mathrm{SR}_0$, no equivalence of $\mathrm{SR}_0$ with unrestricted
uniqueness, and no checker for $\mathrm{SR}_0$.
\end{question}

\begin{question}[A stable introduction at no cost]
Is there an introduction whose mode is preserved and reflected by
substitution, as in \cref{claus:thm:ksubst}, with a complete checker no
larger than the witness-mode checker? The diagonal checker of
\pkg{clause\_guards} is larger by 100 bytes. Is there a guard that
satisfies the conditions of \cref{claus:thm:guard} and is cheaper to test
than the diagonal shape?
\end{question}

\begin{question}[Tree-preserving substitution at no cost]
Does another syntax or certificate design turn the structural property of
$\mathsf{PD}$ (\cref{claus:thm:pdsubst}) into a complete checker below the
current minima? The checker of $\mathsf{PD}$ in
\pkg{476\_byte\_presentation} is smaller than its matched four-family
control but larger than the inherited checker of $\mathsf{ED}$ and than the
compact checkers of other packages, and the package claims no size record
for it.
\end{question}

\begin{question}[One-body contexts without the cost of the gate]
Can the antecedent field be removed, as in
\cref{claus:thm:onebody,claus:thm:homog}, by a representation of
assumptions that needs neither an origin tag nor a gate with lifting, so
that a complete checker becomes smaller? The measured checkers of
\pkg{one\_body} and of the homogeneous variant of \pkg{prefix\_clauses}
are larger than their controls, and a variant of \pkg{one\_body} with
primitive deep reflexivity is larger still (3208 bytes).
\end{question}

\begin{question}[The identity case of the discharge]
Can the identity case of $\varrho$ be replaced by a single arithmetic
pattern that is always active, without a supplied template and without
enlarging the checker or the arithmetic assumptions? The partial operation
$\varrho^{+}$ of \cref{claus:sec:pd} keeps the identity case only for an
$E$ with root $D$; it does not remove the case.
\end{question}

\begin{question}[Combinations not measured]
\pkg{prefix\_clauses} and \pkg{one\_body} ask for the diagonal
introduction on the prefix representation and on one-body contexts, with
a matched measurement. Neither combination is stated or measured in the
packages.
\end{question}

\begin{question}[Sizes of derivations]
\Cref{claus:thm:twocase} bounds inference nodes, not serialized size.
\pkg{raw\_clausal} reports that its translation from the six-rule raw
calculus to $\mathsf{C}_3$ produced at most nine times as many nodes on
every executed example; no general bound is proved. No package bounds the
serialized size of a translated derivation, the cost of closing and
reopening a closure with an unused witness, or the duplication of the
major premise in \cref{claus:thm:fusion} (the translation of
\rulename{Elim} in \cref{claus:thm:pd} avoids that duplication by the
derivation of $\Theta$). \pkg{rootfree\_clausal},
\pkg{reduced\_clause\_calculus} and \pkg{evaluation\_discharge} ask
whether certificates with sharing recover these costs once reference
validation is charged.
\end{question}

The following statements of the packages are negative or unproved and are
recorded as such.
\begin{itemize}
\item No package claims a smallest presentation or a lower bound over
presentations, and none of the checkers of \cref{claus:tab:sizes} is
smaller than the compact codecs of other packages (516 and 508 normalized
bytes under their own input contracts, 788, 780 and 778 bytes for
validated public-goal checkers). The packages ask for a fixed common
contract, comprising goal grammar, scope checking, normalization and
measurement, before any such comparison.
\item The formal statements not present in the Lean development are: the
calculus $\mathsf{C}^{0}_{\mathrm{b}}$ on clauses
(\cref{claus:thm:fixedzero}), the guard criterion and the mirror guard on
clauses (\cref{claus:thm:guard,claus:prop:mirror}), the support quotient
(\cref{claus:prop:support}), the homogeneous calculus
(\cref{claus:thm:homog}), the zero-normal fusion
(\cref{claus:thm:fusionzero}), the propositions
\cref{claus:prop:threenodes,claus:prop:rows}, the obstruction
\cref{claus:thm:reflobs}, the theorem \cref{claus:thm:rigidity}, the
covariance \cref{claus:lem:covar} and the tree-preserving substitution
\cref{claus:thm:pdsubst} (\pkg{476\_byte\_presentation} claims no Lean
or Rocq verification of its results or of its proof compilers;
\cref{claus:thm:ku,claus:thm:pd} are formalized in this repository, see
\cref{claus:rem:pdlean}), and the
correspondence of each Python checker with its calculus.
\item \pkg{reduced\_clause\_calculus} implements the translation of
certificates from raw to zero-normal syntax only; the converse compiler
is not implemented, and the converse is the theorem
\cref{claus:thm:zero}.
\item \pkg{clause\_contexts} reports that neither the ranked
representation nor the support quotient reduced its checker, and
\pkg{prefix\_clauses} that the prefix representation does not reduce
every coordinate against its nested control.
\item \pkg{reduced\_clause\_calculus} asks for a deletion that removes an
arithmetic obligation, a term constructor or a logical operation while
preserving both composites; the deletions of this chapter concern the
proof interface and leave the arithmetic basis unchanged.
\Cref{claus:thm:reflobs,claus:prop:rows} show that, with the other rules
fixed, reflexivity cannot be replaced by finitely many closed seeds and
neither evaluation row of $\mathsf{C}$ can be dropped.
\end{itemize}
